\documentclass[11pt,a4paper,oneside]{report}
% Preambulo del manual del sistema.
\usepackage[utf8]{inputenc}
\usepackage[T1]{fontenc}
\usepackage{lmodern}
\usepackage[spanish,es-noshorthands,es-tabla]{babel}
\usepackage[a4paper,margin=2.4cm]{geometry}
\usepackage{amsmath,amssymb}
\usepackage{booktabs,tabularx,array,multirow,longtable}
\usepackage[table]{xcolor}
\usepackage{tikz}
\usetikzlibrary{arrows.meta,positioning,fit,backgrounds,calc,shapes.geometric,shapes.misc,
                decorations.pathreplacing,patterns,babel,matrix}
\usepackage{pgfplots}
\pgfplotsset{compat=1.18}
\usepackage{listings}
\usepackage[most]{tcolorbox}
\usepackage{enumitem}
\usepackage{caption}
\usepackage{subcaption}
\usepackage{algorithm}
\usepackage{algpseudocode}
\usepackage[section]{placeins}
\usepackage{hyperref}
\hypersetup{colorlinks=true,linkcolor=blue!50!black,urlcolor=blue!50!black,citecolor=blue!50!black}

\setlist{itemsep=2pt,topsep=4pt}
\captionsetup{font=small,labelfont=bf}
\renewcommand{\arraystretch}{1.15}
\newcolumntype{L}{>{\raggedright\arraybackslash}X}

% ---------------------------------------------------------------- colores
\definecolor{datos}{HTML}{4C78A8}
\definecolor{modelo}{HTML}{F58518}
\definecolor{plan}{HTML}{54A24B}
\definecolor{batch}{HTML}{B279A2}
\definecolor{sim}{HTML}{E45756}
\definecolor{exp}{HTML}{72B7B2}
\definecolor{gris}{HTML}{8C8C8C}
\definecolor{m1}{HTML}{4C78A8}
\definecolor{m2}{HTML}{F58518}
\definecolor{m3}{HTML}{54A24B}
\definecolor{m4}{HTML}{E45756}
\definecolor{m5}{HTML}{B279A2}
\definecolor{m6}{HTML}{EECA3B}
\definecolor{m7}{HTML}{9D755D}

% ---------------------------------------------------------------- estilos TikZ
\tikzset{
  caja/.style={draw, rounded corners=2pt, align=center, minimum height=8mm, inner sep=4pt, font=\small},
  flecha/.style={-{Stealth[length=2.2mm]}, thick},
  rotulo/.style={font=\scriptsize, align=center},
  pasillo/.style={fill=gray!8, draw=gray!50},
  bay/.style={draw=gray!60, fill=white, minimum width=4mm, minimum height=3.2mm, inner sep=0pt},
  dock/.style={draw, fill=black!80, text=white, font=\scriptsize\bfseries, minimum width=9mm, minimum height=5mm},
}

% ---------------------------------------------------------------- recuadros
\newtcolorbox{idea}[1][]{colback=plan!6, colframe=plan!70!black, fonttitle=\bfseries\small,
  title={Idea}, #1, boxrule=0.6pt, arc=2pt, left=4pt, right=4pt, top=2pt, bottom=2pt}
\newtcolorbox{supuesto}[1][]{colback=sim!5, colframe=sim!70!black, fonttitle=\bfseries\small,
  title={Supuesto}, #1, boxrule=0.6pt, arc=2pt, left=4pt, right=4pt, top=2pt, bottom=2pt}
\newtcolorbox{hiper}[1]{colback=modelo!6, colframe=modelo!75!black, fonttitle=\bfseries\small,
  title={Hiperparámetros: #1}, boxrule=0.6pt, arc=2pt, left=4pt, right=4pt, top=2pt, bottom=2pt}
\newtcolorbox{codigo}[1]{colback=datos!4, colframe=datos!70!black, fonttitle=\bfseries\small,
  title={En el código: #1}, boxrule=0.6pt, arc=2pt, left=4pt, right=4pt, top=2pt, bottom=2pt}

% ---------------------------------------------------------------- codigo
\lstdefinestyle{py}{
  language=Python, basicstyle=\ttfamily\footnotesize, keywordstyle=\color{datos!80!black}\bfseries,
  commentstyle=\color{plan!60!black}\itshape, stringstyle=\color{modelo!80!black},
  showstringspaces=false, frame=single, rulecolor=\color{gray!40}, backgroundcolor=\color{gray!4},
  breaklines=true, columns=fullflexible, keepspaces=true, numbers=left, numberstyle=\tiny\color{gray},
  xleftmargin=1.5em, framexleftmargin=1.2em, aboveskip=6pt, belowskip=6pt,
  morekeywords={self,None,True,False}
}
\lstset{style=py}
\newcommand{\code}[1]{\texttt{\small #1}}

% ---------------------------------------------------------------- algoritmos en castellano
\floatname{algorithm}{Algoritmo}
\renewcommand{\listalgorithmname}{Lista de algoritmos}
\algrenewcommand\algorithmicrequire{\textbf{Entrada:}}
\algrenewcommand\algorithmicensure{\textbf{Salida:}}
\algrenewcommand\algorithmicfor{\textbf{para}}
\algrenewcommand\algorithmicforall{\textbf{para cada}}
\algrenewcommand\algorithmicdo{\textbf{hacer}}
\algrenewcommand\algorithmicwhile{\textbf{mientras}}
\algrenewcommand\algorithmicif{\textbf{si}}
\algrenewcommand\algorithmicthen{\textbf{entonces}}
\algrenewcommand\algorithmicelse{\textbf{si no}}
\algrenewcommand\algorithmicreturn{\textbf{devolver}}
\algrenewcommand\algorithmicend{\textbf{fin}}
\algrenewcommand\algorithmicrepeat{\textbf{repetir}}
\algrenewcommand\algorithmicuntil{\textbf{hasta}}

% ---------------------------------------------------------------- notacion
\newcommand{\RS}{R_S}
\newcommand{\Aisles}{\mathcal{A}}


\title{\textbf{Slotting basado en afinidad}\\[4pt]
  \large Manual del sistema: datos, modelo, métodos, batching, simulación y experimentos}
\author{Santiago Fada}
\date{Octubre de 2026}

\begin{document}
\maketitle

\begin{abstract}
Este documento explica cómo funciona el sistema que acompaña la tesis: qué datos usa, cómo
los convierte en un problema de optimización, qué métodos resuelven ese problema, cómo se
arman los viajes de picking, cómo se simula el depósito para medir cada propuesta y cómo se
organizan los experimentos. Cada pieza se presenta con su definición, sus hiperparámetros y
un diagrama. El código aparece como mapa de módulos en el capítulo~\ref{cap:arquitectura} y
como fragmentos comentados en el apéndice~\ref{ap:codigo}. Los resultados se resumen en el
capítulo~\ref{cap:resultados}, solo como contexto: el detalle está en los notebooks.
\end{abstract}

\tableofcontents

\chapter{Introducción}
\label{cap:introduccion}

\section{El problema}

En un depósito de e-commerce, un \emph{picker} recorre los pasillos con un carro y junta los
productos de varias órdenes en un mismo viaje. La mayor parte del costo de preparar un pedido
es ese recorrido (Bartholdi y Hackman, 2014). Dos decisiones lo determinan:

\begin{itemize}
  \item \textbf{Slotting}: en qué ubicación del depósito se guarda cada producto (SKU).
  \item \textbf{Batching}: qué órdenes se recogen juntas en un mismo viaje.
\end{itemize}

La tesis estudia el slotting basado en \emph{afinidad}: ubicar cerca los productos que se
piden juntos, además de ubicar cerca del despacho los que se piden mucho. El sistema
descrito aquí permite construir el problema a partir de los datos de ShipHero, resolverlo con
distintos métodos, armar los viajes de picking de distintas formas y medir cada combinación
con una simulación del depósito.

\section{El recorrido completo}

La figura~\ref{fig:pipeline} muestra el camino de los datos hasta una medición. Cada caja es
una etapa del sistema y un capítulo de este documento.

\begin{figure}[ht]
\centering
\begin{tikzpicture}[e/.style={caja, minimum width=29mm, minimum height=11mm}]
  \node[e, fill=datos!15, draw=datos] (tablas) at (0,0) {Cinco tablas\\\scriptsize del depósito};
  \node[e, fill=datos!15, draw=datos] (split) at (3.8,0) {Corte\\\scriptsize train / test};
  \node[e, fill=modelo!15, draw=modelo] (inst) at (7.6,0) {Instancia\\\scriptsize $f$, $c$, $A$, $D$ (de train)};
  \node[e, fill=plan!15, draw=plan] (plan) at (11.4,0) {Plan de slotting\\\scriptsize SKU $\to$ ubicación};
  \node[e, fill=batch!15, draw=batch] (batch) at (11.4,-2.4) {Batching\\\scriptsize órdenes de test $\to$ viajes};
  \node[e, fill=sim!15, draw=sim] (sim) at (7.6,-2.4) {Simulación\\\scriptsize del período de test};
  \node[e, fill=sim!15, draw=sim] (met) at (3.8,-2.4) {Métricas\\\scriptsize caminata por viaje};
  \node[e, fill=exp!20, draw=exp] (expe) at (0,-2.4) {Experimentos\\\scriptsize grillas de configuraciones};
  \draw[flecha] (tablas) -- (split);
  \draw[flecha] (split) -- (inst);
  \draw[flecha] (inst) -- node[above, rotulo]{método} (plan);
  \draw[flecha] (plan) -- node[right, rotulo]{pasillo de\\cada SKU} (batch);
  \draw[flecha] (batch) -- node[above, rotulo]{viajes} (sim);
  \draw[flecha] (plan.south west) -- node[above, sloped, rotulo]{estado inicial} (sim.north east);
  \draw[flecha] (sim) -- (met);
  \draw[flecha] (met) -- (expe);
\end{tikzpicture}
\caption{Del dato a la medición. Los colores se repiten en el resto del documento: azul para
datos, naranja para el modelo, verde para el slotting, violeta para el batching, rojo para la
simulación y celeste para los experimentos.}
\label{fig:pipeline}
\end{figure}

En palabras:

\begin{enumerate}
  \item Se leen las cinco tablas de un depósito y se cortan en un período de entrenamiento
        (\emph{train}) y uno de prueba (\emph{test}) (capítulo~\ref{cap:datos}).
  \item Con train se construye la \emph{instancia}: cuánto se pide cada producto ($f$), cuánto
        cuesta llegar a cada ubicación ($c$), qué productos aparecen juntos ($A$) y qué
        distancia separa a dos ubicaciones ($D$) (capítulo~\ref{cap:modelo}).
  \item Un \emph{método de slotting} usa la instancia para producir un \emph{plan}: la
        ubicación de cada producto (capítulo~\ref{cap:metodos}).
  \item Las órdenes de test se agrupan en viajes, por orden de llegada o con un algoritmo que
        junta órdenes que comparten pasillos (capítulo~\ref{cap:batching}).
  \item La simulación aplica el plan al depósito, reproduce el período de test y mide cuánto
        camina el picker en cada viaje, con una política de recorrido
        (capítulo~\ref{cap:ruteo}) y una política de reposición (capítulo~\ref{cap:simulacion}).
  \item Los experimentos repiten todo lo anterior para muchas combinaciones de datos, métodos
        y parámetros (capítulo~\ref{cap:experimentos}).
\end{enumerate}

\section{Dos medidas distintas}

El sistema separa dos cosas que conviene no confundir:

\begin{itemize}
  \item El \textbf{objetivo} que optimiza un método, calculado sobre train. Es una fórmula
        rápida de evaluar que guía la búsqueda (sección~\ref{sec:formulacion}).
  \item La \textbf{medida} con que se juzga un plan: la caminata simulada sobre los viajes de
        test. Es lo que se reporta.
\end{itemize}

Un método puede mejorar su objetivo y empeorar la medida. Por eso toda comparación del
documento se hace con la simulación.

\section{Cómo leer este documento}

Los capítulos~\ref{cap:datos} a~\ref{cap:simulacion} explican cada etapa: qué es, cómo
funciona, qué hiperparámetros tiene y qué efecto tienen. Cada uno usa tres tipos de
recuadros:

\begin{idea}
La razón de una decisión de diseño.
\end{idea}
\begin{supuesto}
Algo que el sistema asume porque los datos no lo informan. Se nombra siempre.
\end{supuesto}
\begin{hiper}{nombre de la pieza}
Los parámetros de la pieza, sus valores posibles y su valor por defecto.
\end{hiper}

El capítulo~\ref{cap:arquitectura} es el mapa del código, el~\ref{cap:experimentos}
describe cómo se corren los experimentos, y el~\ref{cap:resultados} resume lo que se obtuvo.
El apéndice~\ref{ap:codigo} reúne los fragmentos de código que implementan las ideas
centrales, para quien quiera retomar el proyecto.

% =============================================================================
\chapter{El depósito y los datos}
\label{cap:datos}

Los datos son depósitos sintéticos que ShipHero generó a partir de su operación. Hay dos
versiones. La v1 (abril de 2026) es un único depósito con 30 días de historia y sin
identificador de orden. La v2 (septiembre de 2026) trae tres tamaños, \code{small},
\code{medium} y \code{big}, con seis meses de historia cada uno y el identificador de la orden
en cada línea de picking. En octubre de 2026 se sumó \code{medium-alpha-1.3}, una variante de
\code{medium}. Este documento trabaja con la v2.

% ----------------------------------------------------------------------------- layout
\newcommand{\layout}{%
  \fill[gray!8] (0.3,-0.8) rectangle (8.1,-0.2);
  \fill[gray!8] (0.3,4.75) rectangle (8.1,5.35);
  \fill[gray!8] (0.3,1.98) rectangle (8.1,2.53);
  \foreach \a in {0,...,3} {
    \pgfmathsetmacro{\cx}{2.0*\a+1.2}
    \fill[gray!8] (\cx-0.3,-0.8) rectangle (\cx+0.3,5.35);
    \foreach \j in {0,...,9} {
      \pgfmathsetmacro{\yy}{0.4*\j + (\j>4)*0.55}
      \draw[gray!60, fill=white] (\cx-0.7,\yy) rectangle (\cx-0.32,\yy+0.36);
      \draw[gray!60, fill=white] (\cx+0.32,\yy) rectangle (\cx+0.7,\yy+0.36);
    }
    \pgfmathtruncatemacro{\num}{\a+1}
    \node[font=\scriptsize] at (\cx,5.6) {pasillo \num};
  }
  \node[dock] (dock) at (-0.4,-0.5) {dock};
}

\section{El espacio}

El depósito es una zona de pasillos paralelos (figura~\ref{fig:layout}). A cada lado de un
pasillo hay una hilera de \textbf{bays}, que son columnas de estantería. Un bay se nombra por
su pasillo y su número: \code{A03-07} es el bay 7 del pasillo 3. Los bays impares quedan del
lado de menor coordenada $x$ y los pares enfrente, de modo que \code{A03-07} y \code{A03-08}
están a la misma altura y a distancia cero entre sí.

Cada bay tiene varios \textbf{estantes} y cada estante seis \textbf{bins}. La combinación de
bay, estante y bin es una \textbf{ubicación}, el hueco donde vive un producto
(figura~\ref{fig:bay}). Los \textbf{pasillos transversales} cortan las estanterías y
permiten pasar de un pasillo a otro sin ir hasta la punta. El \textbf{dock} es la mesa de
empaque: ahí empieza y termina cada recorrido.

\begin{figure}[ht]
\centering
\begin{tikzpicture}[scale=0.95]
  \layout
  \node[rotulo, anchor=west] at (8.2,-0.5) {frente};
  \node[rotulo, anchor=west] at (8.2,5.05) {fondo};
  \node[rotulo, anchor=west] at (8.2,2.25) {pasillo\\transversal};
  \draw[datos, very thick] (0.5,0.8) rectangle (0.88,1.16);
  \draw[datos, -{Stealth}] (-0.6,1.6) node[left, rotulo, text=datos]{bay \code{A01-03}\\lado impar} -- (0.5,1.05);
  \draw[datos, very thick] (1.52,0.8) rectangle (1.9,1.16);
  \draw[datos, -{Stealth}] (2.6,-1.3) node[below, rotulo, text=datos]{bay \code{A01-04}, enfrente} -- (1.75,0.8);
  \draw[|-|, gris] (10.3,-0.2) -- node[right, rotulo]{largo $L$\\del pasillo} (10.3,4.75);
\end{tikzpicture}
\caption{Plano esquemático de \code{small}: 4 pasillos, 10 bays por lado y un pasillo
transversal. \code{medium} tiene 8 pasillos y \code{big} 16, con más bays y más
transversales.}
\label{fig:layout}
\end{figure}

\begin{figure}[ht]
\centering
\begin{tikzpicture}[scale=0.55]
  \foreach \s [count=\i from 0] in {A,B,C,D} {
    \node[font=\scriptsize] at (-0.7,\i+0.5) {\s};
    \foreach \b in {1,...,6} {
      \draw[gray!70, fill=white] (\b-1,\i) rectangle (\b,\i+1);
    }
  }
  \foreach \b in {1,...,6} { \node[font=\scriptsize] at (\b-0.5,-0.5) {\b}; }
  \fill[plan!40] (1,1) rectangle (2,2);
  \draw[plan!70!black, very thick] (1,1) rectangle (2,2);
  \draw[-{Stealth}, plan!70!black] (8,1.5) node[right, rotulo, text=black]{ubicación \code{A03-07-B-02}\\estante B, bin 2} -- (2,1.5);
  \node[rotulo] at (3,-1.4) {bins};
  \node[rotulo, rotate=90] at (-1.6,2) {estantes};
\end{tikzpicture}
\caption{Un bay visto de frente: estantes por bins. Cada casilla es una ubicación.}
\label{fig:bay}
\end{figure}

\begin{codigo}{coordenadas y distancias}
La tabla \code{coordinates} da para cada bay su pasillo, número, lado y coordenadas $(x, y)$
en pulgadas. La tabla \code{distances} da la distancia caminando entre cada par de bays,
incluido el dock, calculada como camino mínimo por los pasillos. La distancia de una
ubicación es la de su bay: el movimiento dentro de un bay no se modela.
\end{codigo}

% ----------------------------------------------------------------------------- actores
\section{Productos, órdenes y viajes}

\begin{itemize}
  \item Un \textbf{SKU} es un producto distinto. Al empezar la historia cada SKU ocupa una
        ubicación y cada ubicación tiene a lo sumo un SKU; el 10\% de las ubicaciones está
        vacía.
  \item Un \textbf{merchant} es el vendedor dueño del SKU. El depósito es compartido por
        varios merchants.
  \item Una \textbf{orden} es lo que pidió un cliente: entre 2 y 5 SKUs distintos de un mismo
        merchant, con 1 a 3 unidades de cada uno.
  \item Un \textbf{batch} es un viaje de picking: 25 órdenes que un picker recoge en una sola
        pasada. En los datos son 25 órdenes consecutivas, en el orden en que llegaron, así que
        un batch mezcla merchants y junta entre 60 y 90 SKUs distintos
        (figura~\ref{fig:ordenes}).
  \item Un \textbf{picker} recorre el depósito recogiendo; un \textbf{repositor} repone stock.
\end{itemize}

\begin{figure}[ht]
\centering
\begin{tikzpicture}[x=5mm,y=5mm]
  \foreach \o/\m/\n in {0/m1/3, 1/m3/2, 2/m2/5, 3/m1/4, 4/m4/3} {
    \node[font=\scriptsize, anchor=east] at (-0.4,-\o*1.3) {orden \the\numexpr\o+1\relax};
    \foreach \k in {1,...,\n} { \fill[\m] (\k-1,-\o*1.3-0.45) rectangle (\k-0.1,-\o*1.3+0.45); }
  }
  \node[font=\scriptsize] at (2,-6.9) {$\vdots$ \; (25 órdenes)};
  \draw[decorate, decoration={brace, amplitude=5pt, mirror}] (5.6,-7.5) -- (5.6,0.6);
  \node[caja, fill=batch!12, draw=batch, anchor=west] (b) at (7,-3.2) {batch = un viaje\\\scriptsize la unión de los SKUs\\\scriptsize de sus 25 órdenes};
  \foreach \c/\t [count=\i] in {m1/merchant A, m2/merchant B, m3/merchant C, m4/merchant D} {
    \fill[\c] (15,-\i*0.9+0.3) rectangle (15.7,-\i*0.9+0.9);
    \node[font=\scriptsize, anchor=west] at (15.9,-\i*0.9+0.6) {\t};
  }
\end{tikzpicture}
\caption{Cada orden tiene SKUs de un solo merchant (color). Un batch junta 25 órdenes y, por
lo tanto, SKUs de varios merchants.}
\label{fig:ordenes}
\end{figure}

\begin{idea}
Que dos productos se \emph{compren} juntos es un hecho de la orden. Que se \emph{recojan}
juntos es un hecho del batch, y depende de cómo se arme el batch. Con el identificador de
orden se puede medir lo primero y construir lo segundo con la regla que se quiera
(capítulo~\ref{cap:batching}).
\end{idea}

% ----------------------------------------------------------------------------- tablas
\section{Las cinco tablas}

Cada depósito es una carpeta con cinco tablas (tabla~\ref{tab:tablas}). Se unen por la
ubicación y por el bay (figura~\ref{fig:er}).

\begin{table}[ht]
\centering\small
\begin{tabularx}{\textwidth}{@{}l l L@{}}
\toprule
tabla & una fila por & columnas principales \\
\midrule
\code{coordinates} & bay, y el dock & \code{bay\_id}, \code{aisle}, \code{bay\_number}, \code{side}, \code{x}, \code{y} \\
\code{distances} & par de bays & \code{bay\_a}, \code{bay\_b}, \code{distance\_in} \\
\code{initial\_stock} & ubicación & \code{location\_id}, \code{bay\_id}, \code{sku}, \code{merchant\_account\_id}, \code{units} \\
\code{picking\_events} & línea de pick & \code{batch\_id}, \code{order\_id}, \code{timestamp}, \code{location\_id}, \code{sku}, \code{quantity} \\
\code{replenishment\_events} & reposición & \code{timestamp}, \code{source\_location\_id}, \code{target\_location\_id}, \code{sku}, \code{quantity} \\
\bottomrule
\end{tabularx}
\caption{Las tablas de cada depósito, en formato parquet.}
\label{tab:tablas}
\end{table}

\begin{figure}[ht]
\centering
\begin{tikzpicture}[tabla/.style={draw=datos, fill=datos!8, rounded corners=2pt, align=left, font=\scriptsize, inner sep=4pt}]
  \node[tabla] (co) at (0,0) {\textbf{coordinates}\\bay\_id\\aisle, bay\_number, side\\x, y};
  \node[tabla] (di) at (0,-2.6) {\textbf{distances}\\bay\_a, bay\_b\\distance\_in};
  \node[tabla] (st) at (5,0) {\textbf{initial\_stock}\\location\_id\\bay\_id\\sku, merchant, units};
  \node[tabla] (pe) at (10.5,1.3) {\textbf{picking\_events}\\batch\_id, order\_id\\timestamp\\location\_id, sku, quantity};
  \node[tabla] (re) at (10.5,-1.6) {\textbf{replenishment\_events}\\timestamp\\source / target\_location\_id\\sku, quantity};
  \draw[thick, datos] (st.west) -- node[above, rotulo]{bay\_id} (co.east);
  \draw[thick, datos] (co.south) -- node[right, rotulo]{bay\_id = bay\_a o bay\_b} (di.north);
  \draw[thick, datos] (st.east) -- node[above, sloped, rotulo]{location\_id} (pe.west);
  \draw[thick, datos] (st.east) -- node[below, sloped, rotulo]{source / target} (re.west);
\end{tikzpicture}
\caption{Cómo se unen las tablas.}
\label{fig:er}
\end{figure}

% ----------------------------------------------------------------------------- tamanos
\section{Los cuatro depósitos}

Los cuatro depósitos de la v2 siguen la misma lógica y cambian de escala
(tabla~\ref{tab:tamanos}).

\begin{table}[ht]
\centering\small
\begin{tabular}{@{}lrrrr@{}}
\toprule
 & \code{small} & \code{medium} & \code{medium-alpha-1.3} & \code{big} \\
\midrule
pasillos & 4 & 8 & 8 & 16 \\
bays por lado de pasillo & 10 & 16 & 16 & 24 \\
pasillos transversales & 1 & 2 & 2 & 3 \\
ubicaciones & 1.920 & 7.680 & 7.680 & 27.648 \\
SKUs & 1.728 & 6.912 & 6.912 & 24.883 \\
merchants & 4 & 50 & 50 & 100 \\
pickers & 5 & 20 & 20 & 60 \\
órdenes & 27.150 & 203.625 & 203.625 & 452.500 \\
batches & 1.086 & 8.145 & 8.145 & 18.100 \\
batches por día & --- & 45 & 45 & 100 \\
SKUs que explican el 50\% de las líneas & --- & 318 & 227 & 755 \\
reposiciones & 9.417 & 76.966 & 77.016 & 171.333 \\
\bottomrule
\end{tabular}
\caption{Tamaños de la v2. \code{medium-alpha-1.3} es \code{medium} generado con la demanda
más concentrada en pocos SKUs (parámetro \code{pareto\_alpha} del generador).}
\label{tab:tamanos}
\end{table}

% ----------------------------------------------------------------------------- reposiciones
\section{Stock y reposiciones}

Cada línea de picking descuenta unidades de la ubicación del SKU. Cuando un pick dejaría el
stock en negativo, justo antes ocurre una reposición (figura~\ref{fig:reposicion}). Alrededor
del 45\% de las reposiciones son \emph{mudanzas}: el SKU cambia de ubicación. Por eso el
plano del depósito cambia durante los seis meses.

\begin{figure}[ht]
\centering
\begin{tikzpicture}[node distance=5mm and 8mm]
  \node[caja] (pick) {línea de pick\\del SKU $s$};
  \node[caja, diamond, aspect=2.2, inner sep=1pt, right=of pick] (alc) {¿alcanza\\el stock?};
  \node[caja, right=of alc] (desc) {se descuentan\\las unidades};
  \node[caja, diamond, aspect=2.2, inner sep=1pt, below=of alc] (quedan) {¿quedaban\\unidades?};
  \node[caja, left=of quedan] (rec) {se recarga\\en el lugar};
  \node[caja, fill=sim!10, draw=sim, right=of quedan] (mud) {con probabilidad 0,9 se muda a\\una ubicación libre al azar;\\si no, se recarga en el lugar};
  \draw[flecha] (pick) -- (alc);
  \draw[flecha] (alc) -- node[above, rotulo]{sí} (desc);
  \draw[flecha] (alc) -- node[right, rotulo]{no: reposición} (quedan);
  \draw[flecha] (quedan) -- node[above, rotulo]{sí} (rec);
  \draw[flecha] (quedan) -- node[above, rotulo]{no} (mud);
\end{tikzpicture}
\caption{Cómo generó los datos el simulador de ShipHero. En la simulación de este sistema el
momento, el SKU y la cantidad de cada reposición se copian de los datos; lo único que se
decide es el destino (sección~\ref{sec:reposicion}).}
\label{fig:reposicion}
\end{figure}

% ----------------------------------------------------------------------------- hechos
\section{Hechos verificados de los datos}
\label{sec:hechos}

Varias decisiones del sistema se apoyan en propiedades de los datos. Cada una está
verificada en un notebook de \code{notebooks/data2/}.

\begin{table}[ht]
\centering\small
\begin{tabularx}{\textwidth}{@{}L l@{}}
\toprule
hecho & dónde se verifica \\
\midrule
Cada orden tiene SKUs de un solo merchant (203.625 de 203.625 en \code{medium}). & 04, hecho 1 \\
Todos los batches tienen 25 órdenes. Sale uno cada 32 minutos en \code{medium} y cada 14,4 en \code{big}. & 04, hecho 2; 05, sec.~5 \\
Agrupar las órdenes de a 25 por hora de su primera línea reproduce el 100\% de los batches. & 04, sec.~3.1 \\
Ninguna orden pide el mismo SKU en dos líneas. & 05, sec.~6 \\
En el corte train/test cada SKU está en exactamente una ubicación. & 04, hecho 3 \\
El volumen es constante mes a mes y la popularidad de cada SKU es estable. & 01 \\
La co-ocurrencia por orden no se distingue de la que darían órdenes armadas al azar según la popularidad de cada SKU dentro de su merchant: para los pares con al menos 20 co-ocurrencias esperadas, el cociente entre observado y esperado queda entre 0,65 y 1,19 (percentiles 5 y 99). & 01 \\
\bottomrule
\end{tabularx}
\caption{Propiedades de la v2 que usa el sistema.}
\label{tab:hechos}
\end{table}

\begin{idea}
El último hecho dice que los datos no tienen \emph{afinidad plantada}: dos SKUs aparecen
juntos tanto como lo predice su popularidad. La única estructura conjunta es la del merchant:
dos SKUs de merchants distintos nunca comparten una orden. Esto explica buena parte de los
resultados del capítulo~\ref{cap:resultados}.
\end{idea}

\begin{supuesto}
Los datos no traen la hora en que llega cada orden, solo el orden en que entraron a los
batches. Donde hace falta un tiempo de llegada (la espera de las órdenes en el
capítulo~\ref{cap:batching}) se usa la posición en esa secuencia.
\end{supuesto}

% ----------------------------------------------------------------------------- split
\section{Entrenamiento y prueba}
\label{sec:split}

Los batches se ordenan por la hora de su primera línea y el último 40\% queda como test,
72 días, unas 10 semanas (figura~\ref{fig:split}). Los notebooks de \code{old\_base/} usaron el último
20\%; se pasó a 40\% para ver cómo evoluciona el depósito mientras el plan entra por las
mudanzas (capítulo~\ref{cap:simulacion}). El corte se hace por batch, nunca a mitad de uno, porque el batch es
la unidad que se mide; y como una orden nunca se reparte entre batches, tampoco parte órdenes.
Todo lo que un método aprende (la demanda, la afinidad) sale de train. Todo lo que se reporta
sale de test.

\begin{figure}[ht]
\centering
\begin{tikzpicture}
  \fill[datos!20] (0,0) rectangle (7.2,0.6);
  \fill[sim!25] (7.2,0) rectangle (12,0.6);
  \draw (0,0) rectangle (12,0.6);
  \node[font=\small] at (3.6,0.3) {train: 60\% de los batches};
  \node[font=\small] at (9.6,0.3) {test: 40\%};
  \draw[thick] (7.2,-0.2) -- (7.2,0.9) node[above, rotulo]{corte};
  \node[rotulo, anchor=north west] at (0,-0.1) {1 de enero};
  \node[rotulo, anchor=north east] at (12,-0.1) {30 de junio};
  \draw[flecha, datos] (3.6,-0.6) -- (3.6,-1.2) node[below, rotulo]{demanda $f$, afinidad $A$, estado al corte};
  \draw[flecha, sim] (9.6,-0.6) -- (9.6,-1.2) node[below, rotulo]{caminata simulada};
\end{tikzpicture}
\caption{Corte temporal por batch.}
\label{fig:split}
\end{figure}

Las reposiciones necesitan un corte propio. Un corte por fecha deja ubicaciones con stock
negativo en train, porque algunas mudanzas quedan registradas después de los picks que ya las
presuponen. La función \code{split\_replenishments} corrige el borde: cada ubicación que
terminaría train con stock negativo recupera su primera reposición de test, y se repite hasta
que todas cierran.

\begin{hiper}{corte train/test}
\begin{itemize}
  \item \code{test\_size}: fracción de batches que va a test. Valor por defecto: 0,4 (era 0,2 en
        \code{old\_base/}). Es un campo de \code{InstanceSpec}.
\end{itemize}
\end{hiper}

% =============================================================================
\chapter{El modelo de slotting}
\label{cap:modelo}

Este capítulo convierte los datos de train en un problema de optimización: la
\emph{instancia}. La instancia guarda cuatro parámetros (figura~\ref{fig:parametros}) y
define el objetivo que optimizan los métodos del capítulo~\ref{cap:metodos}.

\begin{figure}[ht]
\centering
\begin{tikzpicture}
  \tikzset{d/.style={caja, fill=datos!12, draw=datos, minimum width=30mm}, p/.style={caja, fill=modelo!15, draw=modelo, minimum width=34mm}}
  \node[d] (pe) at (0,1.6) {\code{picking\_events} de train};
  \node[d] (is) at (0,0.4) {\code{initial\_stock}};
  \node[d] (di) at (0,-0.8) {\code{distances}};
  \node[d] (co) at (0,-2.0) {\code{coordinates}};
  \node[p] (f) at (5.2,2.0) {demanda $f_i$};
  \node[p] (a) at (5.2,1.0) {afinidad $a_{ij}$};
  \node[p] (u) at (5.2,0.0) {SKUs y ubicaciones};
  \node[p] (c) at (5.2,-1.0) {costo de acceso $c_\ell$};
  \node[p] (d) at (5.2,-2.0) {distancia $d_{\ell k}$};
  \node[caja, fill=modelo!25, draw=modelo, very thick, minimum height=14mm] (inst) at (10,0) {\textbf{instancia}\\\code{SlottingInstance}};
  \draw[flecha] (pe.east) -- (f.west);
  \draw[flecha] (pe.east) -- (a.west);
  \draw[flecha] (is.east) -- (u.west);
  \draw[flecha] (is.east) -- (c.west);
  \draw[flecha] (di.east) -- (c.west);
  \draw[flecha] (di.east) -- (d.west);
  \draw[flecha, dashed] (co.east) -- node[below, rotulo, pos=0.4]{si $D$ = \code{aisle}} (d.west);
  \foreach \n in {f,a,u,c,d} { \draw[flecha] (\n.east) -- (inst.west); }
\end{tikzpicture}
\caption{De las tablas a la instancia. Todo sale de train salvo la geometría, que es fija.}
\label{fig:parametros}
\end{figure}

\section{Notación}

\begin{table}[ht]
\centering\small
\begin{tabular}{@{}ll@{}}
\toprule
símbolo & significado \\
\midrule
$I$, $n = |I|$ & SKUs a ubicar (todos los del stock inicial) \\
$\mathcal{L}$, $m = |\mathcal{L}|$ & ubicaciones, $m \ge n$ \\
$f_i$ & demanda del SKU $i$ \\
$c_\ell$ & costo de acceso de la ubicación $\ell$: distancia caminando de su bay al dock \\
$a_{ij}$ & afinidad entre los SKUs $i$ y $j$; simétrica, $a_{ii} = 0$ \\
$d_{\ell k}$ & distancia entre las ubicaciones $\ell$ y $k$, a través de sus bays \\
$x_{i\ell} \in \{0,1\}$ & vale 1 si el SKU $i$ va a la ubicación $\ell$ \\
$\lambda \in [0,1]$ & peso del término de demanda frente al de afinidad \\
\bottomrule
\end{tabular}
\caption{Notación del modelo.}
\end{table}

% ----------------------------------------------------------------------------- f
\section{Demanda $f$}

La demanda dice cuánto se pide cada SKU en train. Hay cuatro formas de contarla, y la
instancia usa una.

\begin{table}[ht]
\centering\small
\begin{tabular}{@{}ll@{}}
\toprule
\code{demand\_metric} & cuenta, para cada SKU \\
\midrule
\code{pick\_lines} & las líneas de picking (por defecto) \\
\code{unique\_orders} & las órdenes distintas que lo piden \\
\code{unique\_batches} & los batches distintos que lo visitan \\
\code{total\_units} & las unidades \\
\bottomrule
\end{tabular}
\caption{Definiciones de la demanda.}
\end{table}

En la v2 \code{unique\_orders} es idéntica a \code{pick\_lines}, porque ninguna orden pide el
mismo SKU en dos líneas (sección~\ref{sec:hechos}).

% ----------------------------------------------------------------------------- c
\section{Costo de acceso $c$}

El costo de una ubicación es la distancia caminando desde su bay hasta el dock, tomada de la
tabla \code{distances}. Todas las ubicaciones de un mismo bay tienen el mismo costo.

% ----------------------------------------------------------------------------- A
\section{Afinidad $A$}
\label{sec:afinidad}

La afinidad mide qué tanto aparecen juntos dos SKUs. Se construye en cuatro pasos
(figura~\ref{fig:afinidad}).

\begin{figure}[ht]
\centering
\begin{tikzpicture}[node distance=5mm and 6mm]
  \node[caja, fill=datos!12, draw=datos] (p) {líneas de\\picking de train};
  \node[caja, fill=modelo!12, draw=modelo, right=of p] (u) {1. unidad\\\scriptsize orden o batch};
  \node[caja, fill=modelo!12, draw=modelo, right=of u] (n) {2. conteos\\\scriptsize $n_{ij}$, $s_i$};
  \node[caja, fill=modelo!12, draw=modelo, right=of n] (m) {3. métrica\\\scriptsize jaccard, cosine,\\\scriptsize cooccurrence};
  \node[caja, fill=modelo!12, draw=modelo, right=of m] (fl) {4. filtro\\\scriptsize top-$k$, mutuo,\\\scriptsize umbral};
  \node[caja, fill=modelo!25, draw=modelo, right=of fl] (a) {$A$};
  \draw[flecha] (p) -- (u); \draw[flecha] (u) -- (n); \draw[flecha] (n) -- (m);
  \draw[flecha] (m) -- (fl); \draw[flecha] (fl) -- (a);
\end{tikzpicture}
\caption{Construcción de la afinidad.}
\label{fig:afinidad}
\end{figure}

\paragraph{1. Unidad de co-ocurrencia.} Dos SKUs co-ocurren cuando aparecen en la misma
\emph{orden} o en el mismo \emph{batch}. Con la orden se mide qué se compra junto. Con el
batch se mide qué se recoge junto, que en los datos mezcla 25 órdenes de varios merchants.

\paragraph{2. Conteos.} Con la matriz binaria $B$ de unidades por SKUs ($B_{ui} = 1$ si el SKU
$i$ está en la unidad $u$), $n_{ij} = (B^\top B)_{ij}$ cuenta las unidades que contienen a los
dos SKUs y el soporte $s_i = (B^\top B)_{ii}$ las que contienen a $i$.

\paragraph{3. Métrica.} Convierte los conteos en un puntaje:
\[
\text{cooccurrence: } a_{ij} = n_{ij}, \qquad
\text{jaccard: } a_{ij} = \frac{n_{ij}}{s_i + s_j - n_{ij}}, \qquad
\text{cosine: } a_{ij} = \frac{n_{ij}}{\sqrt{s_i\, s_j}}.
\]
El conteo crudo favorece a los SKUs populares, que co-ocurren mucho solo por ser populares.
Jaccard y coseno lo corrigen dividiendo por los soportes.

\paragraph{4. Filtro.} La matriz completa es densa y casi todo es ruido de baja frecuencia.
Los filtros conservan los vínculos fuertes y devuelven siempre una matriz simétrica:
\begin{itemize}
  \item \code{top\_k}: cada SKU conserva sus $k$ vecinos de mayor afinidad, y un vínculo
        sobrevive si está en el top-$k$ de \emph{alguno} de los dos (unión).
  \item \code{mutual\_top\_k}: sobrevive solo si está en el top-$k$ de \emph{los dos}
        (intersección).
  \item \code{threshold}: sobrevive si supera un valor mínimo.
\end{itemize}

\paragraph{Un ejemplo.} Cuatro órdenes: $\{a,b\}$, $\{a,b,c\}$, $\{a,c\}$ y $\{b,d\}$. Los
soportes son $s_a = 3$, $s_b = 3$, $s_c = 2$, $s_d = 1$, y los pares que co-ocurren son
$n_{ab} = 2$, $n_{ac} = 2$, $n_{bc} = 1$, $n_{bd} = 1$. La figura~\ref{fig:filtros} muestra
la afinidad Jaccard y lo que conserva cada filtro con $k = 1$.

\begin{figure}[ht]
\centering
\begin{tikzpicture}[sku/.style={circle, draw, fill=white, minimum size=6mm, font=\small}]
  \begin{scope}
    \node[sku] (a) at (0,1.4) {$a$}; \node[sku] (b) at (1.8,1.4) {$b$};
    \node[sku] (c) at (0,0) {$c$}; \node[sku] (d) at (1.8,0) {$d$};
    \draw (a) -- node[above, rotulo]{0,50} (b);
    \draw (a) -- node[left, rotulo]{0,67} (c);
    \draw (b) -- node[above, rotulo, sloped]{0,25} (c);
    \draw (b) -- node[right, rotulo]{0,33} (d);
    \node[rotulo] at (0.9,-0.8) {Jaccard, sin filtro};
  \end{scope}
  \begin{scope}[xshift=4.2cm]
    \node[sku] (a) at (0,1.4) {$a$}; \node[sku] (b) at (1.8,1.4) {$b$};
    \node[sku] (c) at (0,0) {$c$}; \node[sku] (d) at (1.8,0) {$d$};
    \draw[very thick, modelo] (a) -- (b); \draw[very thick, modelo] (a) -- (c); \draw[very thick, modelo] (b) -- (d);
    \node[rotulo] at (0.9,-0.8) {\code{top\_k}, $k = 1$};
  \end{scope}
  \begin{scope}[xshift=8.4cm]
    \node[sku] (a) at (0,1.4) {$a$}; \node[sku] (b) at (1.8,1.4) {$b$};
    \node[sku] (c) at (0,0) {$c$}; \node[sku] (d) at (1.8,0) {$d$};
    \draw[very thick, modelo] (a) -- (c);
    \node[rotulo] at (0.9,-0.8) {\code{mutual\_top\_k}, $k = 1$};
  \end{scope}
\end{tikzpicture}
\caption{El mejor vecino de $a$ es $c$; el de $b$, $a$; el de $c$, $a$; el de $d$, $b$. La
unión conserva $ac$, $ab$ y $bd$. La intersección solo $ac$, el único par que se elige
mutuamente.}
\label{fig:filtros}
\end{figure}

\begin{hiper}{afinidad}
\begin{itemize}
  \item \code{cooccurrence\_col}: \code{"order\_id"} o \code{"batch\_id"}. Por defecto
        \code{"batch\_id"}, la única posible en la v1; con la v2 se usa \code{"order\_id"}.
  \item \code{affinity\_metric}: \code{"jaccard"} (por defecto), \code{"cosine"} o
        \code{"cooccurrence"}.
  \item \code{affinity\_filter}: \code{TopKFilter(k)}, \code{MutualTopKFilter(k)},
        \code{ThresholdFilter(t)} o ninguno (matriz completa, por defecto). Valor base de
        los experimentos: \code{TopKFilter(10)}.
  \item \code{normalize}: si $f$ y $A$ se pasan a proporciones (sección~\ref{sec:escala}).
        Por defecto \code{False}; los experimentos usan \code{True}.
\end{itemize}
\end{hiper}

% ----------------------------------------------------------------------------- D
\section{Distancia entre productos $D$}
\label{sec:distanciaD}

El término de afinidad cobra la distancia entre las ubicaciones de dos SKUs afines. Hay dos
formas de medirla (figura~\ref{fig:distD}):

\begin{itemize}
  \item \code{walking}: la distancia caminando entre los dos bays, de la tabla
        \code{distances}.
  \item \code{aisle}: lo que cuesta ese par en un recorrido en S (capítulo~\ref{cap:ruteo}),
        donde cada pasillo con algún pick se recorre entero. Dos SKUs en el mismo pasillo no se
        agregan nada; en pasillos distintos cuestan un pasillo entero más el cruce lateral:
        \[
        d(\ell, k) = \begin{cases}
          0 & \text{si } \ell \text{ y } k \text{ están en el mismo pasillo,} \\
          L + |x_{\text{pasillo}(\ell)} - x_{\text{pasillo}(k)}| & \text{si no,}
        \end{cases}
        \]
        con $L$ el largo de un pasillo y $x$ la coordenada lateral del centro del pasillo.
\end{itemize}

\begin{figure}[ht]
\centering
\begin{tikzpicture}[scale=0.8]
  \foreach \cx in {0,2.4} {
    \fill[gray!8] (\cx-0.35,0) rectangle (\cx+0.35,4);
    \foreach \j in {0,...,7} {
      \draw[gray!60, fill=white] (\cx-0.8,\j*0.5) rectangle (\cx-0.37,\j*0.5+0.45);
      \draw[gray!60, fill=white] (\cx+0.37,\j*0.5) rectangle (\cx+0.8,\j*0.5+0.45);
    }
  }
  \fill[modelo] (-0.8,0.5) rectangle (-0.37,0.95); \node[rotulo] at (-1.15,0.72) {$i$};
  \fill[modelo] (0.37,3.0) rectangle (0.8,3.45); \node[rotulo] at (1.15,3.22) {$j$};
  \fill[plan] (2.77,2.0) rectangle (3.2,2.45); \node[rotulo] at (3.5,2.22) {$k$};
  \node[rotulo, align=left, anchor=west] at (4.3,3.0) {$i, j$ en el mismo pasillo:\\
     \code{walking} = el tramo entre sus bays\\ \code{aisle} = 0};
  \node[rotulo, align=left, anchor=west] at (4.3,1.2) {$i, k$ en pasillos distintos:\\
     \code{walking} = camino más corto\\ \code{aisle} = $L$ + distancia lateral};
\end{tikzpicture}
\caption{Las dos distancias entre productos.}
\label{fig:distD}
\end{figure}

\begin{hiper}{distancia entre productos}
\begin{itemize}
  \item \code{pair\_distance\_model}: \code{"walking"} (por defecto) o \code{"aisle"}. Los
        modelos distintos de \code{walking} necesitan la tabla \code{coordinates}. Solo lo usa
        el término de afinidad; el costo de acceso $c$ y la simulación usan siempre la
        geometría física.
\end{itemize}
\end{hiper}

% ----------------------------------------------------------------------------- formulacion
\section{La formulación}
\label{sec:formulacion}

El plan es una asignación de cada SKU a una ubicación. El objetivo combina dos términos:

\begin{equation}
\min_{x} \; C(x) \;=\; \lambda \underbrace{\sum_{i \in I} \sum_{\ell \in \mathcal{L}} f_i\, c_\ell\, x_{i\ell}}_{L:\ \text{demanda por costo de acceso}}
\;+\; (1-\lambda) \underbrace{\sum_{i,j \in I} \sum_{\ell, k \in \mathcal{L}} a_{ij}\, d_{\ell k}\, x_{i\ell}\, x_{jk}}_{Q:\ \text{afinidad por distancia entre ubicaciones}}
\label{eq:qap}
\end{equation}
\begin{align*}
\text{s.a.}\quad & \sum_{\ell \in \mathcal{L}} x_{i\ell} = 1 \quad \forall i \in I && \text{(cada SKU en una ubicación)}\\
& \sum_{i \in I} x_{i\ell} \le 1 \quad \forall \ell \in \mathcal{L} && \text{(a lo sumo un SKU por ubicación)}\\
& x_{i\ell} \in \{0, 1\}.
\end{align*}

$L$ es chico cuando los SKUs más pedidos están cerca del dock. $Q$ es chico cuando los pares
afines están cerca entre sí. Los dos tiran en direcciones distintas, y $\lambda$ decide cuánto
pesa cada uno: con $\lambda = 1$ solo importa la demanda; con $\lambda = 0$, solo la
afinidad.

El problema es un \emph{Quadratic Assignment Problem} (QAP) en la forma de Koopmans y
Beckmann (1957), con término lineal. Es NP-difícil, y el límite práctico de la resolución
exacta ronda las 30 instalaciones (Loiola et al., 2007). Kofler et al. (2014) usan la misma
combinación de popularidad y afinidad, ponderada por un parámetro; Li et al. (2015) deciden a
nivel de fila con capacidad, como los métodos por pasillo de la
sección~\ref{sec:aislefill}. Con 6.912 SKUs y 7.680 ubicaciones (\code{medium}) tiene unos 53 millones de
variables binarias. El método exacto del sistema (sección~\ref{sec:exacto}) prueba
optimalidad hasta unos 20 SKUs. Por eso los métodos que se usan en la práctica son
heurísticas.

% ----------------------------------------------------------------------------- escala
\section{La escala de los dos términos y el significado de $\lambda$}
\label{sec:escala}

$L$ suma líneas por pulgadas y $Q$ suma afinidades (con Jaccard, números menores que 1) por
pulgadas. Sin corregir, $L$ es unas 1.500 veces más grande que $Q$, y $\lambda = 0{,}5$ no
pesa la afinidad en absoluto.

La corrección (\code{normalize=True}) divide $f$ y $A$ por su suma:
\[
\tilde f_i = \frac{f_i}{\sum_k f_k}, \qquad \tilde a_{ij} = \frac{a_{ij}}{\sum_{k,l} a_{kl}}.
\]
Con eso $L$ es la distancia media al dock de una línea de picking y $Q$ la distancia media
entre dos SKUs de un par afín, las dos en pulgadas. $\lambda = 0{,}5$ significa que las dos
distancias medias pesan igual.

% ----------------------------------------------------------------------------- deltas
\section{Cambios incrementales}
\label{sec:deltas}

Las heurísticas mueven un SKU por vez y necesitan saber cuánto cambia $C$ sin recalcularlo
entero. Sea $\ell_i$ la ubicación del SKU $i$ y $N(i)$ sus vecinos de afinidad.

\paragraph{Intercambiar dos SKUs $a$ y $b$.}
\[
\Delta L = (f_a - f_b)(c_{\ell_b} - c_{\ell_a}), \qquad
\Delta Q = 2 \sum_{k \neq a,b} (a_{ak} - a_{bk})\,(d_{\ell_b \ell_k} - d_{\ell_a \ell_k}).
\]

\paragraph{Mudar un SKU $s$ a una ubicación libre $y$.}
\[
\Delta L = f_s\,(c_y - c_{\ell_s}), \qquad
\Delta Q = 2 \sum_{k \in N(s)} a_{sk}\,(d_{y \ell_k} - d_{\ell_s \ell_k}).
\]

En los dos casos $\Delta C = \lambda\, \Delta L + (1-\lambda)\, \Delta Q$. El factor 2 cuenta
cada par en sus dos órdenes y vale porque $A$ es simétrica. Con un filtro top-$k$ cada SKU
tiene pocos vecinos, así que cada cálculo cuesta $O(k)$ en lugar de $O(n)$.

\begin{idea}
$C$ es un \emph{sustituto}: se calcula con datos de train y es rápido de evaluar. El plan se
juzga con otra medida, la caminata simulada sobre test. Si las dos medidas no van en la misma
dirección, la simulación manda.
\end{idea}

% =============================================================================
\chapter{El recorrido del picker}
\label{cap:ruteo}

Para medir un plan hay que saber por dónde camina el picker en cada viaje. Eso lo decide una
\emph{política de ruteo}: recibe los bays que el batch tiene que visitar y devuelve la
distancia caminada desde el dock y de vuelta. El sistema tiene dos.

\section{Recorrido en S estricto}
\label{sec:sshape}

Es el recorrido de la operación de ShipHero. El picker sale del dock y va al frente del
pasillo con picks que está más a la izquierda. Lo recorre entero hasta el fondo, cruza al
siguiente pasillo con picks y lo recorre entero hacia el frente, y así alternando. Los
pasillos sin picks se saltean. Si recorrió una cantidad impar de pasillos termina en el
fondo, y vuelve al frente recorriendo el último otra vez. Después regresa al dock
(figura~\ref{fig:sshape}).

\begin{figure}[ht]
\centering
\begin{tikzpicture}[scale=0.95]
  \layout
  % picks del batch
  \foreach \x/\y in {0.51/0.4, 1.53/3.75, 2.51/1.2, 3.53/4.15, 6.51/0.8, 7.53/3.35} {
    \fill[sim] (\x,\y) rectangle (\x+0.37,\y+0.36);
  }
  % recorrido
  \draw[sim, very thick, -{Stealth}] (0.05,-0.5) -- (1.2,-0.5) -- (1.2,5.05) -- (3.2,5.05) -- (3.2,-0.5)
        -- (7.2,-0.5) -- (7.2,5.05);
  \draw[sim, very thick, dashed, -{Stealth}] (7.32,5.05) -- (7.32,-0.62) -- (0.05,-0.62);
  \node[rotulo, text=sim, anchor=west] at (8.2,0.8) {vuelve por el último\\(3 pasillos: impar)};
  \node[rotulo, anchor=west] at (8.2,3.8) {el pasillo 3 no tiene\\picks: se saltea};
\end{tikzpicture}
\caption{Recorrido en S para un batch con picks (rojo) en los pasillos 1, 2 y 4.}
\label{fig:sshape}
\end{figure}

Para un batch que visita los pasillos $S = \{a_1 < \dots < a_m\}$:

\begin{equation}
\RS(S) = \underbrace{d_0(a_1)}_{\text{ida}} + \underbrace{m\,L}_{\text{recorrer cada pasillo}}
+ \underbrace{x_{a_m} - x_{a_1}}_{\text{cruces}}
+ \underbrace{[\,m \text{ impar}\,]\,L}_{\text{volver al frente}} + \underbrace{d_0(a_m)}_{\text{regreso}}
\label{eq:rs}
\end{equation}
\[
d_0(a) = |x_a - x_{\text{dock}}| + (y_0 - y_{\text{dock}}),
\]

con $L$ el largo de un pasillo, $x_a$ la coordenada lateral de su centro, $y_0$ la del frente
de los pasillos y $d_0(a)$ la distancia del dock al frente del pasillo $a$. Los cruces suman
$x_{a_m} - x_{a_1}$ porque los pasillos se visitan de izquierda a derecha.

\begin{idea}
Ningún término de~\eqref{eq:rs} depende de en qué bay del pasillo está cada pick. \textbf{El
largo del recorrido depende solo de qué pasillos se visitan.} Para el slotting, eso cambia el
objetivo: no importa poner un producto cerca del frente, importa que los batches entren a pocos
pasillos.
\end{idea}

La S es una heurística de ruteo, no el recorrido más corto. Valle et al. (2017) miden que
las mejores heurísticas quedan en promedio a 2,4\% del ruteo óptimo, y que la brecha baja
cuando los pedidos son grandes, como los batches de 60 a 90 SKUs de estos datos.

\paragraph{El efecto de la paridad.} Con una cantidad impar de pasillos hay que recorrer el
último dos veces. Por eso agregar un pasillo puede costar casi nada, o nada. La
tabla~\ref{tab:paridad} lo muestra con la geometría de \code{medium}.

\begin{table}[ht]
\centering\small
\begin{tabular}{@{}lrl@{}}
\toprule
pasillos visitados & $\RS$ (pulgadas) & \\
\midrule
$\{1\}$ & 4.480 & uno solo: se recorre dos veces \\
$\{1, 2\}$ & 4.864 & solo 384 más que un pasillo \\
$\{1, 3, 5\}$ & 9.568 & \\
$\{1, 3, 4, 5\}$ & 9.568 & agregar el 4 no cuesta nada \\
los 8 pasillos & 17.824 & \\
\bottomrule
\end{tabular}
\caption{Valores de $\RS$ en \code{medium} ($L = 1.776$ pulgadas). Calculados en el notebook 04.}
\label{tab:paridad}
\end{table}

\paragraph{Cálculo.} Como $\RS$ depende solo del conjunto de pasillos, se calcula una vez para
todos los conjuntos posibles y se guarda en una tabla. Cada conjunto se codifica como un entero
con un bit por pasillo: $\{1, 3\}$ son los bits 0 y 2, o sea $101_2 = 5$. Con 8 pasillos son
$2^8 = 256$ valores; con 16, $65.536$. La unión de dos conjuntos es el \emph{o} bit a bit de
sus enteros. Esto es lo que hace rápido al algoritmo de batching del
capítulo~\ref{cap:batching}.

\section{Recorrido serpenteante (\code{snake})}

Es la política con la que se midieron los resultados anteriores a octubre de 2026. Ordena los
bays a visitar por (pasillo, número de bay) y suma el camino más corto entre cada par
consecutivo, más la ida y la vuelta al dock:
\begin{equation}
R_E = D(\text{dock}, \beta_1) + \sum_{t=1}^{T-1} D(\beta_t, \beta_{t+1}) + D(\beta_T, \text{dock}).
\end{equation}
Como el camino más corto entre el último pick de un pasillo y el primero del siguiente suele
volver por el frente, en la práctica el picker entra en cada pasillo hasta el pick más lejano y
sale por donde entró (figura~\ref{fig:snake}). Esta medida cobra la \emph{profundidad} a la que
se entra en cada pasillo, que la S no cobra.

\begin{figure}[ht]
\centering
\begin{tikzpicture}[scale=0.95]
  \layout
  \foreach \x/\y in {0.51/0.4, 1.53/3.75, 2.51/1.2, 3.53/4.15, 6.51/0.8, 7.53/3.35} {
    \fill[sim] (\x,\y) rectangle (\x+0.37,\y+0.36);
  }
  \draw[datos, very thick, -{Stealth}] (0.05,-0.5) -- (1.15,-0.5) -- (1.15,3.93);
  \draw[datos, very thick, -{Stealth}] (1.25,3.93) -- (1.25,-0.45) -- (3.15,-0.45) -- (3.15,4.33);
  \draw[datos, very thick, -{Stealth}] (3.25,4.33) -- (3.25,-0.45) -- (7.15,-0.45) -- (7.15,3.53);
  \draw[datos, very thick, dashed, -{Stealth}] (7.25,3.53) -- (7.25,-0.6) -- (0.05,-0.6);
  \node[rotulo, text=datos, anchor=west] at (8.2,2.0) {entra hasta el pick\\más lejano y vuelve};
\end{tikzpicture}
\caption{El mismo batch con la política \code{snake}: la caminata depende de la profundidad
de cada pick, no solo de los pasillos.}
\label{fig:snake}
\end{figure}

\section{Qué cambia para el slotting}

Con la S, lo que el slotting puede bajar es la cantidad de pasillos que visita un batch. Un
razonamiento sencillo muestra qué conviene.

\begin{supuesto}
Si las $N$ líneas de un batch se eligieran al azar e independientes, cada una con probabilidad
igual a la demanda de su SKU, un pasillo $a$ que concentra una fracción $F_a$ de la demanda
quedaría sin visitar con probabilidad $(1 - F_a)^N$, y la cantidad esperada de pasillos
visitados sería
\[
\mathbb{E}[\text{pasillos}] = \sum_{a \in \Aisles} \bigl(1 - (1 - F_a)^N\bigr).
\]
Los datos son compatibles con ese supuesto (sección~\ref{sec:hechos}).
\end{supuesto}

Cada sumando es una función cóncava de $F_a$. Para una misma demanda total, la suma es menor
cuando la demanda se concentra en pocos pasillos que cuando se reparte. El greedy por demanda
hace lo contrario: pone lo más pedido cerca del dock, que es el frente de \emph{todos} los
pasillos. De ahí salen los planes que llenan el depósito pasillo por pasillo
(sección~\ref{sec:aislefill}). Este es el mismo argumento que Schleyer y Gue (2012) usan para
la distribución de pasillos visitados con recorrido en S.

\begin{hiper}{ruteo}
\begin{itemize}
  \item \code{routing} del simulador: \code{"snake"} (por defecto, para no cambiar los
        resultados anteriores) o \code{"s\_shape"}. Todos los resultados desde el notebook 04
        usan \code{"s\_shape"}.
\end{itemize}
\end{hiper}

% =============================================================================
\chapter{Los métodos de slotting}
\label{cap:metodos}

Un método de slotting recibe la instancia y devuelve un \emph{plan}: la ubicación de cada SKU.
Todos cumplen el mismo contrato, \code{solve(instance) -> Assignment}, y se registran con un
nombre para pedirlos sin acoplarse a su clase (sección~\ref{sec:registros}). La
tabla~\ref{tab:metodos} los resume.

\begin{table}[ht]
\centering\small
\begin{tabularx}{\textwidth}{@{}l L c >{\raggedright\arraybackslash}p{4.4cm}@{}}
\toprule
nombre & idea & afinidad & hiperparámetros \\
\midrule
\code{current} / vigente & lo que hay en el depósito & no & --- \\
\code{random} & cada SKU en una ubicación sorteada: el plan que no usa datos & no & semilla \\
\code{demand\_greedy} & lo más pedido, lo más cerca del dock & no & --- \\
\code{linear\_assignment} & óptimo exacto de $L$ & no & --- \\
\code{aisle\_fill} & lo más pedido, llenando un pasillo por vez & no & --- \\
\code{merchant\_aisle\_fill} & igual, agrupando cada merchant & no & --- \\
\code{bilevel} & agrupar, repartir bloques, resolver cada bloque & según el bloque & agrupamiento, costo de zona, método interno \\
\code{swap\_search} & búsqueda local sobre $C$ & sí & $\lambda$, vecindad, plan inicial, pasadas \\
\code{exact\_qap} & resolver $C$ exacto con un solver & sí & $\lambda$, tiempo, gap \\
\bottomrule
\end{tabularx}
\caption{Métodos de slotting del sistema.}
\label{tab:metodos}
\end{table}

\section{El plan vigente}

Es la línea de base de la operación: dónde está cada SKU. Hay dos versiones. El método
\code{current} lee el stock del día cero. Los experimentos usan el \emph{vigente al corte}: la
ubicación de cada SKU en el estado real del depósito al terminar train, calculado por la
simulación (sección~\ref{sec:estados}). Como en el corte cada SKU está en una sola ubicación
(sección~\ref{sec:hechos}), ese estado define un plan.

\section{Greedy por demanda y asignación lineal}

\code{demand\_greedy} ordena los SKUs por demanda de mayor a menor, ordena las ubicaciones por
costo de acceso de menor a mayor, y manda el $i$-ésimo SKU a la $i$-ésima ubicación. No usa la
afinidad: es el método natural cuando $\lambda = 1$.

\begin{idea}
Con $\lambda = 1$ el objetivo es $L = \sum_i f_i\, c_{\ell_i}$. Por la desigualdad de
reordenamiento, una suma de productos se minimiza emparejando el mayor de un lado con el menor
del otro. El greedy es entonces óptimo para $L$, sin necesidad de un solver.
\end{idea}

\code{linear\_assignment} resuelve el mismo problema con el algoritmo húngaro. Sirve como
control: su valor de $L$ coincide con el del greedy, y en las simulaciones su caminata queda a
0,4\% o menos. Trabaja con una matriz densa de SKUs por ubicaciones, así que en \code{big}
(24.883 por 27.648) no entra en memoria y no se corre.

\section{Llenado por pasillos}
\label{sec:aislefill}

Los dos métodos de esta sección salen del capítulo~\ref{cap:ruteo}: con recorrido en S
conviene concentrar la demanda en pocos pasillos. Ordenan las ubicaciones de otra forma:

\begin{enumerate}
  \item Cada pasillo tiene un \emph{rango}: el costo de su ubicación más barata, o sea qué tan
        cerca del dock está su frente. Con el dock en una esquina, el pasillo 1 es el primero.
  \item Las ubicaciones se ordenan por rango de su pasillo y, dentro del pasillo, por costo.
        Así se llena un pasillo entero antes de pasar al siguiente.
\end{enumerate}

Los dos métodos difieren en cómo ordenan los SKUs:

\begin{itemize}
  \item \code{aisle\_fill}: por demanda de mayor a menor. El pasillo 1 recibe los SKUs más
        pedidos, el 2 los siguientes, y así.
  \item \code{merchant\_aisle\_fill}: por merchant y después por demanda. Los merchants se
        ordenan por demanda total, y dentro de cada merchant los SKUs por demanda. Como cada
        merchant tiene menos SKUs que las ubicaciones de un pasillo, queda entero en un pasillo
        o partido entre dos consecutivos.
\end{itemize}

\begin{idea}
Cada orden es de un solo merchant (sección~\ref{sec:hechos}). Si los SKUs de un merchant
están en el mismo pasillo, una orden se recoge entrando a un solo pasillo. Con FIFO eso
sirve de poco, porque un batch junta 25 órdenes de merchants distintos. Con re-batching
(capítulo~\ref{cap:batching}) es lo que permite armar batches que entran a dos o tres
pasillos.
\end{idea}

\begin{figure}[ht]
\centering
\resizebox{\textwidth}{!}{%
\begin{tikzpicture}[x=1cm,y=1cm]
  \foreach \panel/\titulo [count=\p from 0] in {g/\code{demand\_greedy}, a/\code{aisle\_fill}, m/\code{merchant\_aisle\_fill}, b/\code{bilevel merchant/dock}} {
    \begin{scope}[xshift=\p*4.4cm]
      \foreach \a in {0,...,3} {
        \foreach \s in {0,1} {
          \foreach \j in {0,...,9} {
            \pgfmathsetmacro{\xx}{\a*0.85+\s*0.32}
            \pgfmathsetmacro{\yy}{\j*0.28}
            \pgfmathtruncatemacro{\k}{\a*20+\j*2+\s}
            \pgfmathsetmacro{\dist}{(\xx+0.5)+(\yy+0.5)}
            \ifcase\p
              \pgfmathtruncatemacro{\pct}{max(5,100-(\dist-1)*100/3.6)}
              \fill[plan!\pct] (\xx,\yy) rectangle (\xx+0.3,\yy+0.26);
            \or
              \pgfmathtruncatemacro{\pct}{max(5,100-\k*100/80)}
              \fill[plan!\pct] (\xx,\yy) rectangle (\xx+0.3,\yy+0.26);
            \or
              \pgfmathtruncatemacro{\mm}{int(\k/12)+1}
              \pgfmathtruncatemacro{\pp}{100-mod(\k,12)*55/12}
              \fill[m\mm!\pp] (\xx,\yy) rectangle (\xx+0.3,\yy+0.26);
            \or
              \pgfmathtruncatemacro{\mm}{min(7,int((\dist-1)*7/3.6)+1)}
              \fill[m\mm!75] (\xx,\yy) rectangle (\xx+0.3,\yy+0.26);
            \fi
            \draw[white, line width=0.2pt] (\xx,\yy) rectangle (\xx+0.3,\yy+0.26);
          }
        }
      }
      \node[dock, minimum width=6mm, minimum height=3mm, font=\tiny\bfseries] at (-0.35,-0.3) {dock};
      \node[rotulo] at (1.55,3.15) {\titulo};
    \end{scope}
  }
\end{tikzpicture}}
\caption{Dónde termina la demanda con cada regla, en un depósito de 4 pasillos. Verde más
oscuro: SKUs más pedidos. Colores: merchants. El greedy reparte lo más pedido por el frente de
todos los pasillos; \code{aisle\_fill} lo concentra en el pasillo 1;
\code{merchant\_aisle\_fill} pone cada merchant en un bloque de pasillo; \code{bilevel} con
costo por distancia al dock reparte cada merchant en una franja que cruza todos los pasillos.}
\label{fig:reglas}
\end{figure}

\begin{hiper}{llenado por pasillos}
\begin{itemize}
  \item No tienen hiperparámetros propios. Necesitan la tabla \code{coordinates} para saber
        el pasillo de cada bay. La demanda que usan es la de la instancia
        (\code{demand\_metric}).
\end{itemize}
\end{hiper}

\section{Bi-nivel}
\label{sec:bilevel}

\code{bilevel} descompone el problema en dos niveles (figura~\ref{fig:bilevel}):

\begin{enumerate}
  \item \textbf{Agrupar} los SKUs con una estrategia de agrupamiento:
        \begin{itemize}
          \item \code{merchant}: un grupo por merchant.
          \item \code{demand\_class}: clases A, B y C por demanda acumulada. Se ordenan los
                SKUs por demanda y los cortes están donde la demanda acumulada pasa el 80\% y
                el 95\%.
        \end{itemize}
  \item \textbf{Repartir bloques de ubicaciones} entre los grupos. Con $y_{\ell g} = 1$ si la
        ubicación $\ell$ va al grupo $g$, $F_g$ la demanda del grupo y $n_g$ su cantidad de
        SKUs:
        \[
        \min \sum_{\ell, g} F_g\, \kappa_\ell\, y_{\ell g} \quad \text{s.a.} \quad
        \sum_g y_{\ell g} \le 1 \;\; \forall \ell, \qquad
        \sum_\ell y_{\ell g} = n_g \;\; \forall g, \qquad y \in \{0,1\}.
        \]
        Es un problema de transporte: su matriz es totalmente unimodular y la relajación
        lineal ya da una solución entera. Lo resuelve el solver. $\kappa_\ell$ es el
        \emph{costo de zona}, el criterio con que se ordenan las ubicaciones.
  \item \textbf{Resolver cada bloque} como un problema independiente, restringido a los SKUs
        del grupo y a sus ubicaciones, con cualquier método. En los experimentos, el greedy.
\end{enumerate}

\begin{figure}[ht]
\centering
\begin{tikzpicture}[node distance=5mm and 9mm]
  \node[caja, fill=modelo!15, draw=modelo] (i) {instancia};
  \node[caja, fill=plan!12, draw=plan, right=of i] (g) {1. agrupar\\\scriptsize merchant o A/B/C};
  \node[caja, fill=plan!12, draw=plan, right=of g] (t) {2. repartir bloques\\\scriptsize transporte, con solver};
  \node[caja, fill=plan!12, draw=plan, right=of t, yshift=7mm] (z1) {3. bloque 1\\\scriptsize método interno};
  \node[caja, fill=plan!12, draw=plan, right=of t, yshift=-7mm] (z2) {3. bloque $G$\\\scriptsize método interno};
  \node[caja, fill=plan!25, draw=plan, right=of z1, yshift=-7mm] (p) {plan};
  \draw[flecha] (i) -- (g); \draw[flecha] (g) -- (t);
  \draw[flecha] (t) -- (z1); \draw[flecha] (t) -- (z2);
  \draw[flecha] (z1) -- (p); \draw[flecha] (z2) -- (p);
  \node[rotulo] at ($(z1)!0.5!(z2)$) {$\vdots$};
\end{tikzpicture}
\caption{Los tres pasos de \code{bilevel}.}
\label{fig:bilevel}
\end{figure}

\paragraph{El costo de zona decide la forma de los bloques.} Con $\kappa_\ell$ = distancia al
dock, el grupo de más demanda recibe las ubicaciones más cercanas al dock, que están en el
frente de todos los pasillos: los bloques son franjas (último panel de la
figura~\ref{fig:reglas}). Con $\kappa_\ell$ = posición en el orden de \code{aisle\_fill}, los
bloques son pasillos consecutivos. La combinación \code{bilevel merchant/aisle\_order} es la
versión con solver de \code{merchant\_aisle\_fill}.

\begin{hiper}{bilevel}
\begin{itemize}
  \item \code{clustering}: \code{merchant} o \code{demand\_class}; esta última con
        \code{thresholds}, por defecto $(0{,}8;\ 0{,}95)$.
  \item \code{zone\_solver}: el método de cada bloque; en los experimentos,
        \code{demand\_greedy}.
  \item \code{location\_cost} ($\kappa$): ninguno (distancia al dock, por defecto) o un vector
        por ubicación. En los experimentos se llama \code{zone\_cost} y vale \code{dock} o
        \code{aisle\_order}.
\end{itemize}
\end{hiper}

\section{Búsqueda por intercambios}
\label{sec:swap}

\code{swap\_search} parte de un plan y lo mejora moviendo SKUs, guiada por el objetivo $C$ de
la ecuación~\eqref{eq:qap}. La idea es acercar cada SKU a sus vecinos de afinidad: para un
SKU $i$ y cada vecino $k$, los destinos candidatos son las ubicaciones del bay donde está $k$.
Si un candidato está libre, $i$ se muda; si está ocupado por $j$, $i$ y $j$ se intercambian,
aunque $i$ y $j$ no sean afines. Se aplica el primer cambio que baje $C$
(algoritmo~\ref{alg:swap}).

\begin{algorithm}[ht]
\caption{Búsqueda por intercambios (\code{swap\_search}).}
\label{alg:swap}
\begin{algorithmic}[1]
\Require instancia, plan inicial $x$, $\lambda$, \code{max\_passes}, vecindad
\For{$p = 1, \dots, \code{max\_passes}$}
  \State mejoró $\gets$ falso
  \ForAll{SKU $i$ con vecinos de afinidad}
    \ForAll{vecino $k$ de $i$}
      \State candidatos $\gets$ ubicaciones del bay de $k$ \Comment{vecindad \code{bay}}
      \State \hphantom{candidatos $\gets$} o solo la ubicación de $k$ \Comment{vecindad \code{neighbor}}
      \ForAll{$\ell \in$ candidatos}
        \If{$\ell$ libre \textbf{y} $\Delta C(\text{mudar } i \to \ell) < 0$}
          \State mudar $i$ a $\ell$; mejoró $\gets$ verdadero; pasar al siguiente vecino
        \ElsIf{$\ell$ ocupada por $j \ne i$ \textbf{y} $\Delta C(\text{intercambiar } i, j) < 0$}
          \State intercambiar $i$ y $j$; mejoró $\gets$ verdadero; pasar al siguiente vecino
        \EndIf
      \EndFor
    \EndFor
  \EndFor
  \If{no mejoró} \Return $x$ \EndIf
\EndFor
\State \Return $x$
\end{algorithmic}
\end{algorithm}

Cada $\Delta C$ se calcula con las fórmulas de la sección~\ref{sec:deltas}, en tiempo
proporcional a la cantidad de vecinos. La vecindad \code{neighbor} es la versión original:
solo intercambia $i$ con su vecino $k$, lo que pone a $i$ en el bay de $k$ pero saca a $k$ de
ahí (figura~\ref{fig:vecindad}).

\begin{figure}[ht]
\centering
\begin{tikzpicture}[x=6mm,y=6mm]
  \begin{scope}
    \foreach \j in {0,...,4} { \draw[gray!60] (0,\j) rectangle (1,\j+1); }
    \fill[modelo] (0.1,2.1) rectangle (0.9,2.9); \node[font=\scriptsize, text=white] at (0.5,2.5) {$k$};
    \fill[gray!25] (0.1,0.1) rectangle (0.9,0.9);
    \node[font=\scriptsize] at (0.5,0.5) {$j$};
    \node[draw, circle, minimum size=5mm, font=\scriptsize, fill=plan!20] (i) at (4,2.5) {$i$};
    \draw[flecha, plan] (i) to[bend right=20] (1,0.5);
    \draw[flecha, plan] (i) to[bend left=10] (1,1.5);
    \draw[flecha, plan] (i) to[bend left=20] (1,3.5);
    \draw[flecha, plan] (i) to[bend left=30] (1,4.5);
    \node[rotulo, align=center] at (2,-0.9) {\code{bay}: cualquier ubicación del bay de $k$\\(libre: mudanza; ocupada: intercambio)};
  \end{scope}
  \begin{scope}[xshift=6cm]
    \foreach \j in {0,...,4} { \draw[gray!60] (0,\j) rectangle (1,\j+1); }
    \fill[modelo] (0.1,2.1) rectangle (0.9,2.9); \node[font=\scriptsize, text=white] at (0.5,2.5) {$k$};
    \node[draw, circle, minimum size=5mm, font=\scriptsize, fill=plan!20] (i) at (4,2.5) {$i$};
    \draw[{Stealth}-{Stealth}, thick, plan] (i) -- (1,2.5);
    \node[rotulo, align=center] at (2,-0.9) {\code{neighbor}: solo intercambiar $i$ con $k$};
  \end{scope}
\end{tikzpicture}
\caption{Las dos vecindades de \code{swap\_search} para un SKU $i$ y su vecino $k$.}
\label{fig:vecindad}
\end{figure}

\begin{hiper}{swap\_search}
\begin{itemize}
  \item \code{lam} ($\lambda$): peso de $L$ en $C$; por defecto 0,5. Con $\lambda = 1$ y el
        greedy como plan inicial no cambia nada, porque el greedy ya es óptimo para $L$.
  \item \code{neighborhood}: \code{"bay"} (por defecto) o \code{"neighbor"}.
  \item \code{seed\_method}: el plan inicial; por defecto \code{demand\_greedy}. En los
        experimentos también \code{aisle\_fill} y \code{merchant\_aisle\_fill}.
  \item \code{max\_passes}: tope de pasadas; por defecto 50. Se corta antes si una pasada
        no mejora.
  \item Además dependen de la instancia: la afinidad (métrica, unidad, filtro), la
        normalización y la distancia $D$.
\end{itemize}
El tiempo crece con la cantidad de vecinos: en \code{medium}, con \code{top\_k:10} tarda unos
35 minutos; sin filtro, 7,7 horas.
\end{hiper}

\section{Exacto}
\label{sec:exacto}

\code{exact\_qap} escribe el programa~\eqref{eq:qap} y lo resuelve con un solver de
programación entera (Gurobi). Sirve como vara para medir las heurísticas en instancias
chicas: con $\lambda = 0{,}01$ prueba optimalidad en segundos hasta unos 20 SKUs y corta por
tiempo a los 30; con $\lambda = 0$ no cierra en un minuto ni con 20. En esas
subinstancias la búsqueda por intercambios quedó a menos de 0,2\% del óptimo.

\begin{hiper}{exact\_qap}
\begin{itemize}
  \item \code{lam} ($\lambda$), por defecto 0,5.
  \item \code{time\_limit}: segundos antes de devolver la mejor solución encontrada.
  \item \code{mip\_gap}: brecha relativa aceptada.
  \item La licencia del solver se lee de las variables \code{SOLVER\_LICENSE\_*} del archivo
        \code{.env}, que no se versiona.
\end{itemize}
\end{hiper}

% =============================================================================
\chapter{El batching}
\label{cap:batching}

El batching decide qué órdenes recoge juntas un picker en un mismo viaje. Una política de
batching recibe las líneas de picking de un período y devuelve las mismas líneas con un
\code{batch\_id} nuevo. No cambia nada más: los SKUs, las cantidades y las horas quedan como
estaban, así que el resultado entra a la simulación igual que los datos originales.

\section{FIFO}

Es lo que hacen los datos. Las órdenes se ordenan por la hora de su primera línea y se agrupan
de a 25 consecutivas. La clase \code{FifoBatching} reproduce el 100\% de los batches de los
datos (sección~\ref{sec:hechos}).

\section{Re-batching: qué se decide}

En vez de cerrar un batch apenas hay 25 órdenes, se espera a juntar $W \cdot 25$ órdenes
consecutivas, y recién ahí se decide cómo repartirlas en $W$ batches de 25. Ese grupo es un
\textbf{pool} y $W$ es la \textbf{ventana} (figura~\ref{fig:pools}). Con $W = 1$ el pool es
un solo batch y no hay nada que decidir: es FIFO.

\begin{figure}[ht]
\centering
\begin{tikzpicture}[x=1.15mm,y=1mm]
  % FIFO
  \node[rotulo, anchor=east] at (-2,4) {FIFO};
  \foreach \i in {0,...,99} {
    \pgfmathtruncatemacro{\b}{int(\i/25)+1}
    \fill[m\b!70] (\i,2) rectangle (\i+0.8,6);
  }
  % re-armado
  \node[rotulo, anchor=east] at (-2,-6) {re-armado};
  \foreach \i in {0,...,99} {
    \pgfmathtruncatemacro{\b}{mod(\i*7+3,4)+1}
    \fill[m\b!70] (\i,-8) rectangle (\i+0.8,-4);
  }
  \draw[decorate, decoration={brace, amplitude=4pt}] (0,8) -- (99.8,8) node[midway, above=3pt, rotulo]{un pool: $W \cdot 25$ órdenes consecutivas ($W = 4$)};
  \draw[-{Stealth}] (0,-11) -- (100,-11) node[right, rotulo]{orden de llegada};
  \foreach \b in {1,2,3,4} {
    \fill[m\b!70] (14+\b*14,-16) rectangle (17+\b*14,-14);
    \node[rotulo, anchor=west] at (17.5+\b*14,-15) {batch \b};
  }
\end{tikzpicture}
\caption{Cada barra es una orden y su color, el batch al que va. Con FIFO los batches son
tramos consecutivos. Con re-batching, las 100 órdenes del pool se reparten entre los 4
batches según los pasillos que necesitan.}
\label{fig:pools}
\end{figure}

\paragraph{Qué queda fijo.} Las mismas órdenes con sus mismas líneas; 25 órdenes por batch
(ShipHero arma batches de hasta 25 órdenes, una por tote); la misma cantidad de batches; y el
slotting, del que el re-batching solo lee en qué pasillo está cada SKU.

\section{Por qué acorta la caminata}

Para cada orden $o$, $A(o)$ es el conjunto de pasillos donde están sus SKUs según el plan. Un
batch entra a todos los pasillos que necesita alguna de sus órdenes, o sea a la unión de sus
$A(o)$, y con recorrido en S su caminata depende solo de esa unión
(capítulo~\ref{cap:ruteo}). Con \code{merchant\_aisle\_fill} casi todas las órdenes necesitan
un solo pasillo; con FIFO cada batch junta 25 órdenes cualesquiera y entre todas necesitan
casi todos los pasillos. Juntar órdenes que necesitan los mismos pasillos reduce la unión
(figura~\ref{fig:pool_ejemplo}).

\begin{figure}[ht]
\centering
\begin{tikzpicture}[x=5mm,y=5mm]
  \newcommand{\fila}[3]{%
    \node[rotulo, anchor=east] at (-0.3,#1+0.5) {#2};
    \foreach \a in {1,...,8} { \draw[gray!50] (\a-1,#1) rectangle (\a,#1+1); }
    \foreach \a in {#3} { \fill[batch!70] (\a-1+0.08,#1+0.08) rectangle (\a-0.08,#1+0.92); }
  }
  \begin{scope}
    \node[font=\small\bfseries] at (4,5.6) {FIFO};
    \fila{3.6}{batch 1}{1,2,3,4,5,6,7}
    \fila{2.4}{batch 2}{1,2,3,4,5,6,7,8}
    \fila{1.2}{batch 3}{1,2,3,4,5,6,7,8}
    \fila{0}{batch 4}{1,2,3,4,5,6,7}
    \foreach \a in {1,...,8} { \node[rotulo] at (\a-0.5,-0.6) {\a}; }
    \node[rotulo] at (4,-1.5) {pasillo};
    \node[rotulo, anchor=west] at (8.3,2.3) {total 70.528};
  \end{scope}
  \begin{scope}[xshift=7.5cm]
    \node[font=\small\bfseries] at (4,5.6) {re-armado};
    \fila{3.6}{batch 1}{5,6,7,8}
    \fila{2.4}{batch 2}{3,4,5,6}
    \fila{1.2}{batch 3}{1,5}
    \fila{0}{batch 4}{2,3,4}
    \foreach \a in {1,...,8} { \node[rotulo] at (\a-0.5,-0.6) {\a}; }
    \node[rotulo] at (4,-1.5) {pasillo};
    \node[rotulo, anchor=west] at (8.3,2.3) {total 35.872};
  \end{scope}
\end{tikzpicture}
\caption{El primer pool de test de \code{medium} con $W = 4$ y el plan
\code{merchant\_aisle\_fill} (notebook 04). Casillas llenas: pasillos que visita cada batch. La
caminata total del pool, en pulgadas, baja casi a la mitad.}
\label{fig:pool_ejemplo}
\end{figure}

\section{El problema de optimización}

\paragraph{Datos de un pool.} El conjunto de órdenes $P$ con $|P| = 25\,W$, los batches
$K = \{1, \dots, W\}$ y $A(o)$ para cada orden.

\paragraph{Variables.} $z_{ok} = 1$ si la orden $o$ va al batch $k$, y 0 si no.

\paragraph{Pasillos de un batch.} $M_k = \bigcup_{o\,:\,z_{ok} = 1} A(o)$.

\paragraph{Problema.}
\begin{align*}
\min_{z}\quad & \sum_{k \in K} \RS(M_k) && \text{(caminata total del pool)}\\
\text{s.a.}\quad & \sum_{k \in K} z_{ok} = 1 \quad \forall o \in P && \text{(cada orden en un batch)}\\
& \sum_{o \in P} z_{ok} = 25 \quad \forall k \in K && \text{(batches de 25 órdenes)}\\
& z_{ok} \in \{0,1\}.
\end{align*}
El último pool del período puede tener menos órdenes; ahí el último batch queda más chico,
igual que con FIFO.

\paragraph{Por qué no se resuelve exacto.} La cantidad de formas de repartir el pool es
\[
\frac{(25W)!}{(25!)^W\, W!}:
\]
unas $10^{13}$ con $W = 2$, $10^{55}$ con $W = 4$ y
$10^{452}$ con $W = 16$. Además $\RS(M_k)$ no es lineal en $z$: depende de qué pasillos hay
en la unión, de cuántos son y de su paridad. Por eso se usa una heurística, que no garantiza
el óptimo.

\section{El algoritmo}

\code{AisleCoverBatching} trabaja pool por pool en dos pasos (algoritmo~\ref{alg:batching}).

\paragraph{1. Construcción.} Los batches se arman de a uno. Se abre el batch con la orden libre
que sola pide más caminata, $o^\ast = \arg\max_{o} \RS(A(o))$. Después, con $M$ los pasillos
que ya visita, se agrega la orden libre que menos agrega,
\[
\Delta(o \mid M) = \RS\bigl(M \cup A(o)\bigr) - \RS(M),
\]
hasta tener 25 órdenes; en un empate gana la que llegó antes. Abrir con una semilla y crecer
por menor incremento es la familia de algoritmos \emph{seed} de order batching (Yang et al.,
2020; Deng et al., 2025). Que la semilla sea la orden más cara es una elección: es la que más
restringe con quién conviene juntarla.

\paragraph{2. Intercambios.} Se prueba cambiar una orden $o$ del batch $k_1$ por una orden $o'$
del batch $k_2$. Con $M_{k}^{-o}$ los pasillos de $k$ sin la orden $o$ (la unión de las
demás, porque otra orden puede usar los mismos pasillos), el cambio en la caminata es
\[
\delta(o, o') = \RS\bigl(M_{k_1}^{-o} \cup A(o')\bigr) + \RS\bigl(M_{k_2}^{-o'} \cup A(o)\bigr)
- \RS(M_{k_1}) - \RS(M_{k_2}).
\]
Para cada orden se busca la $o'$ de otro batch con menor $\delta$, y si $\delta < 0$ se
aplica. Se repite hasta una pasada sin intercambios, con un tope. Cada intercambio baja
estrictamente la caminata y los repartos son finitos, así que termina; el resultado es un
óptimo local, no necesariamente el óptimo del problema.

\begin{algorithm}[ht]
\caption{Re-batching de un pool (\code{AisleCoverBatching}).}
\label{alg:batching}
\begin{algorithmic}[1]
\Require órdenes del pool en orden de llegada, $A(o)$ de cada una, tabla de $\RS$, tamaño $Q = 25$
\State libres $\gets$ todas las órdenes
\While{quedan libres} \Comment{1. construcción}
  \State $o^\ast \gets$ la libre de mayor $\RS(A(o))$; abrir batch con $o^\ast$; $M \gets A(o^\ast)$
  \While{el batch tiene menos de $Q$ órdenes \textbf{y} quedan libres}
    \State $o \gets$ la libre de menor $\Delta(o \mid M)$; agregarla; $M \gets M \cup A(o)$
  \EndWhile
\EndWhile
\Repeat \Comment{2. intercambios}
  \ForAll{orden $o$, en el batch $k_1$}
    \State $o' \gets$ la orden de otro batch con menor $\delta(o, o')$
    \If{$\delta(o, o') < 0$} intercambiar $o$ y $o'$ \EndIf
  \EndFor
\Until{una pasada sin intercambios o \code{max\_passes} pasadas}
\end{algorithmic}
\end{algorithm}

La figura~\ref{fig:traza} muestra la construcción del primer batch del ejemplo. Las primeras
once órdenes necesitan los pasillos 7 u 8 y no agregan caminata. En el paso 12 ya no queda
ninguna así, y la más barata abre el pasillo 6: $\Delta = 3.552 = 2L$, porque el batch pasa de
2 a 3 pasillos, uno más para recorrer y, por ser impar, uno más para volver. En el paso 14
entra el pasillo 5 sin costo: con 4 pasillos ya no hay que volver por el último.

\begin{figure}[ht]
\centering
\begin{tikzpicture}
\begin{axis}[width=12cm, height=4.6cm, xlabel={orden agregada (paso)}, ylabel={$\RS$ del batch},
  xmin=0.5, xmax=25.5, ymin=6000, ymax=11800, xtick={1,5,10,12,14,20,25}, ytick={7168,10720},
  scaled y ticks=false, yticklabel style={/pgf/number format/.cd, 1000 sep={.}}, grid=major, grid style={gray!20},
  tick label style={font=\scriptsize}, label style={font=\small}]
  \addplot[const plot, batch, very thick] coordinates {(0.5,7168) (11.5,7168) (11.5,10720) (25.5,10720)};
  \node[rotulo, anchor=south west] at (axis cs:1,7250) {pasillos 7 y 8};
  \node[rotulo, anchor=south west] at (axis cs:12,10800) {pasillo 6 ($\Delta = 2L$)};
  \node[rotulo, anchor=north west] at (axis cs:14.2,10650) {pasillo 5 ($\Delta = 0$)};
\end{axis}
\end{tikzpicture}
\caption{Caminata del primer batch del ejemplo mientras se construye (notebook 04,
sección 3.6).}
\label{fig:traza}
\end{figure}

\section{Cómo se calcula rápido}

El paquete hace las mismas cuentas con tres recursos (figura~\ref{fig:mascaras}):

\begin{itemize}
  \item \textbf{Máscaras.} Cada $A(o)$ es un entero con un bit por pasillo, y la unión es el
        \emph{o} bit a bit.
  \item \textbf{Tabla de $\RS$.} $\RS$ de cualquier conjunto se lee de la tabla calculada una
        vez (sección~\ref{sec:sshape}).
  \item \textbf{Conteos por pasillo.} Para cada batch $k$ y pasillo $a$ se guarda $c_{ka}$, la
        cantidad de órdenes del batch que necesitan ese pasillo. El batch sin la orden $o$
        entra a $M_k^{-o} = \{a : c_{ka} - [a \in A(o)] > 0\}$, sin recorrer sus 25 órdenes.
\end{itemize}

\begin{figure}[ht]
\centering
\begin{tikzpicture}[x=5mm,y=5mm, bit/.style={draw, minimum size=5mm, inner sep=0pt, font=\scriptsize}]
  \foreach \v [count=\i from 0] in {0,0,0,0,0,1,0,1} { \ifnum\v=1 \def\col{batch!40}\else\def\col{white}\fi \node[bit, fill=\col] at (\i,2) {\v}; }
  \node[rotulo, anchor=east] at (-0.8,2) {$A(o) = \{6, 8\}$};
  \foreach \v [count=\i from 0] in {0,0,1,0,0,1,0,0} { \ifnum\v=1 \def\col{batch!40}\else\def\col{white}\fi \node[bit, fill=\col] at (\i,0.8) {\v}; }
  \node[rotulo, anchor=east] at (-0.8,0.8) {$M = \{3, 6\}$};
  \draw (-0.5,0.1) -- (7.5,0.1);
  \foreach \v [count=\i from 0] in {0,0,1,0,0,1,0,1} { \ifnum\v=1 \def\col{batch!40}\else\def\col{white}\fi \node[bit, fill=\col] at (\i,-0.6) {\v}; }
  \node[rotulo, anchor=east] at (-0.8,-0.6) {$M \cup A(o)$: \emph{o} bit a bit};
  \foreach \a in {1,...,8} { \node[rotulo] at (\a-1,3.1) {\a}; }
  \node[rotulo, anchor=west] at (8.2,-0.6) {$\to$ $\RS$ = tabla[entero]};
\end{tikzpicture}
\caption{Conjuntos de pasillos como máscaras de bits (el pasillo 1 a la izquierda).}
\label{fig:mascaras}
\end{figure}

Con eso evaluar un intercambio son cuatro lecturas de la tabla, y se evalúa contra todas las
órdenes del pool a la vez. Re-armar todo el período de test de \code{medium} tarda unos 2
segundos.

\section{Cuánto espera una orden}
\label{sec:espera}

Juntar más órdenes antes de decidir tiene un costo: la orden espera más. Como los datos no
traen la hora de llegada, se fija una regla:

\begin{supuesto}
\textbf{Llegada.} Las órdenes llegan en el orden de la FIFO y parejas, 25 por \emph{franja}.
Una franja es el tiempo entre un batch y el siguiente: 32 minutos en \code{medium}, 14,4 en
\code{big}. \textbf{Salida.} Un pool se completa cuando llega su última orden. Desde ahí sus
$W$ batches salen de a uno por franja, al mismo ritmo que con FIFO, para que los pickers
tengan la misma carga; sale primero el batch con la orden que llegó antes.
\end{supuesto}

Para la orden que llega en el lugar $r$ (contando desde 0), en el pool $p$, en un batch que sale
en el turno $t$ ($t = 0$ el primero), la espera en franjas es
\[
\text{espera} = \underbrace{\frac{r_{\max}(p) + 1}{25}}_{\text{se completa el pool}}
+ \underbrace{t}_{\text{turno de su batch}} - \underbrace{\frac{r}{25}}_{\text{llega}},
\]
con $r_{\max}(p)$ el lugar de la última orden del pool (figura~\ref{fig:espera}). Con FIFO la
espera va de $1/25$ de franja a una franja: la primera orden de un batch espera a las otras 24.

\begin{figure}[ht]
\centering
\begin{tikzpicture}[x=1.6cm,y=1cm]
  \draw[-{Stealth}] (0,0) -- (4.6,0) node[right, rotulo]{tiempo (franjas)};
  \foreach \t in {0,...,4} { \draw (\t,0.08) -- (\t,-0.08) node[below, rotulo]{\t}; }
  \fill[datos!25] (0,0.3) rectangle (2,0.7); \node[rotulo] at (1,0.5) {llegan las 50 órdenes del pool};
  \fill[batch!40] (2,0.9) rectangle (3,1.3); \node[rotulo] at (2.5,1.1) {sale el batch 1};
  \fill[batch!25] (3,1.5) rectangle (4,1.9); \node[rotulo] at (3.5,1.7) {sale el batch 2};
  \draw[flecha, sim] (0,2.4) -- node[above, rotulo]{espera de la orden 1 si va en el batch 1: 2 franjas} (2,2.4);
\end{tikzpicture}
\caption{Ventana $W = 2$. Con FIFO la orden 1 saldría en la franja 1, al llegar la orden 25.}
\label{fig:espera}
\end{figure}

\begin{table}[ht]
\centering\small
\begin{tabular}{@{}lrrrrrr@{}}
\toprule
ventana $W$ & 1 (FIFO) & 2 & 4 & 8 & 16 & 32 \\
\midrule
espera media en \code{medium} (min) & 17 & 49 & 113 & 240 & 496 & 1.007 \\
espera máxima en \code{medium} (min) & 32 & 95 & 220 & 462 & 881 & --- \\
espera media en \code{big} (min) & 7 & 22 & 51 & 108 & 223 & 453 \\
\bottomrule
\end{tabular}
\caption{Espera de las órdenes según la ventana, con la regla de arriba (notebooks 04 y 05).
En \code{big} sale un batch más seguido, así que la misma ventana significa menos espera.}
\label{tab:espera}
\end{table}

\begin{hiper}{batching}
\begin{itemize}
  \item \code{batch\_size} ($Q$): órdenes por batch; 25, como en la operación.
  \item \code{window} ($W$): batches FIFO que se re-arman juntos; 1 equivale a FIFO. Valores
        probados: 1, 2, 4, 8, 16 y 32.
  \item \code{max\_passes}: tope de pasadas de intercambios; por defecto 20; con 0 queda solo
        la construcción.
  \item Entradas: el pasillo de cada SKU según el plan (\code{sku\_aisle}) y la tabla de $\RS$
        (\code{aisle\_set\_cost}). El re-batching agrupa siempre con la caminata en S, también
        cuando después se mide con \code{snake}.
\end{itemize}
\end{hiper}

% =============================================================================
\chapter{La simulación}
\label{cap:simulacion}

La simulación es la medida de todo el sistema. Parte del estado real del depósito al
terminar train y reproduce el período de test: cada línea de picking y cada reposición, en
orden. De cada reposición los datos dan el momento, el SKU y la cantidad; la simulación decide
solo adónde van las unidades. Para cada batch mide la caminata con una política de ruteo.

\textbf{El régimen.} Ningún depósito se reordena entero de un día para el otro, así que el plan
no se aplica de una vez: el depósito arranca como está al corte y el plan entra por las
mudanzas que ocurren de todos modos. Cuando un SKU se agota y hay que darle lugar, la regla
\code{plan\_aisle} lo lleva a su pasillo del plan (sección~\ref{sec:reposicion}). Como el pasillo de cada SKU cambia durante
el test, los batches se arman pool por pool con el pasillo que cada SKU tiene en ese momento
(sección~\ref{sec:rolling}). Aplicar el plan entero en el corte, el régimen de los notebooks de
\code{old\_base/}, queda como el \emph{techo} de cada plan.

\section{Los objetos}

Cuatro objetos pasan de una etapa a la otra (figura~\ref{fig:objetos}). La distinción entre
\code{Assignment} y \code{WarehouseState} es la central: el primero es el plan, sin unidades y
sin dinámica; el segundo es el depósito real, que cambia.

\begin{figure}[ht]
\centering
\begin{tikzpicture}[node distance=8mm and 14mm, every node/.append style={font=\footnotesize}]
  \node[caja, fill=modelo!15, draw=modelo, text width=29mm] (i) {\code{SlottingInstance}\\\scriptsize el problema;\\\scriptsize inmutable};
  \node[caja, fill=plan!15, draw=plan, text width=29mm, right=of i] (a) {\code{Assignment}\\\scriptsize el plan:\\\scriptsize SKU $\to$ ubicación};
  \node[caja, fill=sim!15, draw=sim, text width=29mm, right=of a] (s) {\code{WarehouseState}\\\scriptsize SKU y unidades\\\scriptsize por ubicación; cambia};
  \node[caja, fill=sim!25, draw=sim, text width=29mm, below=of s] (r) {\code{SimulationResult}\\\scriptsize métricas y registros};
  \draw[flecha] (i) -- node[above, rotulo]{método} (a);
  \draw[flecha] (a) -- node[above, rotulo]{\code{reslot\_state}} (s);
  \draw[flecha] (s) -- node[right, rotulo]{\code{Simulator.run}} (r);
\end{tikzpicture}
\caption{Los objetos de la simulación.}
\label{fig:objetos}
\end{figure}

Un SKU puede ocupar varias ubicaciones a la vez (pasa en los datos durante una mudanza) y una
ubicación aloja a lo sumo un SKU.

\section{El estado inicial}
\label{sec:estados}

\begin{table}[ht]
\centering\small
\begin{tabularx}{\textwidth}{@{}l L@{}}
\toprule
función & produce \\
\midrule
\code{state\_from\_initial\_stock} & el estado del día cero, leído del stock inicial \\
\code{state\_after\_history} & el estado al terminar train. No recorre los eventos uno por uno: lo calcula por aritmética (stock inicial más entregas menos picks, por ubicación), y asigna cada ubicación al último SKU que se repuso ahí. Así no le afectan los timestamps desordenados de algunas mudanzas \\
\code{reslot\_state} & el estado que resulta de aplicar un plan: cada SKU se muda a su ubicación del plan con las unidades que tenía. El costo de esa mudanza no se cuenta \\
\bottomrule
\end{tabularx}
\caption{Cómo se arma el estado inicial de una simulación.}
\end{table}

En el régimen vigente la simulación arranca de \code{state\_after\_history}: el depósito real al
corte. \code{reslot\_state} se usa solo para el techo; en los notebooks de \code{old\_base/}
era el punto de partida de todas las corridas.

\section{Los eventos}

Las tablas se convierten en una secuencia ordenada por tiempo con dos tipos de evento
(figura~\ref{fig:eventos}):

\begin{itemize}
  \item \code{PickBatch}: un batch completo, con sus líneas (SKU, cantidad, hora). Se procesa
        de una vez y su caminata se mide al final.
  \item \code{ReplenEvent}: una reposición (hora, SKU, cantidad). Solo falta decidir su
        destino.
\end{itemize}

Dentro de un batch, antes de cada línea se aplican todas las reposiciones con hora anterior o
igual. El generador de los datos inserta cada reposición justo antes de la línea que la
dispara; intercalarlas más grueso, por ejemplo al inicio del batch, encontraría al SKU con
stock y convertiría mudanzas en recargas.

\begin{figure}[ht]
\centering
\begin{tikzpicture}[x=1.1cm]
  \draw[-{Stealth}] (0,0) -- (11,0) node[right, rotulo]{tiempo};
  \foreach \x in {0.5,1.2,1.9,2.6,3.3} { \fill[sim] (\x,0.2) rectangle (\x+0.3,0.5); }
  \foreach \x in {6,6.7,7.4,8.1,8.8} { \fill[sim] (\x,0.2) rectangle (\x+0.3,0.5); }
  \draw[decorate, decoration={brace, amplitude=4pt}] (0.5,0.7) -- (3.6,0.7) node[midway, above=4pt, rotulo]{batch 1: líneas de pick};
  \draw[decorate, decoration={brace, amplitude=4pt}] (6,0.7) -- (9.1,0.7) node[midway, above=4pt, rotulo]{batch 2};
  \foreach \x in {2.3, 7.2} { \draw[thick, plan, -{Stealth}] (\x,-0.7) -- (\x,-0.05); }
  \node[rotulo, text=plan] at (2.3,-1.0) {reposición};
  \node[rotulo, text=plan] at (7.2,-1.0) {reposición};
  \node[rotulo, align=left, anchor=west] at (3.9,-0.6) {se aplica justo antes\\de la línea siguiente};
\end{tikzpicture}
\caption{Picks y reposiciones intercalados a su hora exacta.}
\label{fig:eventos}
\end{figure}

\section{La regla de picking}

Cuando un SKU ocupa varias ubicaciones, el picker va a la de menor stock que cubra la línea
sola. Si ninguna la cubre, junta de varias, de menor a mayor stock. Lo que el stock no cubre
se cuenta como \emph{faltante}, sin detener la simulación. Una ubicación que queda en cero se
libera si el SKU tiene otra; la única ubicación de un SKU no se libera.

\section{Las políticas de reposición}
\label{sec:reposicion}

La política decide adónde van las unidades de cada reposición. Es el único punto donde un
plan de slotting actúa durante el test.

\begin{table}[ht]
\centering\small
\begin{tabularx}{\textwidth}{@{}l L L@{}}
\toprule
política & qué hace & representa \\
\midrule
\code{plan\_aisle} & si al SKU le quedan unidades, recarga en el lugar; si se agotó, lo lleva a una ubicación libre de su pasillo del plan, la más cercana a su ubicación del plan, o al pasillo más cercano con lugar & un sistema que implementa el plan de a una mudanza; es la regla del régimen vigente \\
\code{home} & a la ubicación del SKU en el plan; si está ocupada por otro, a la libre más cercana a ella & un sistema que conoce el plan y lo mantiene; se usa para el techo \\
\code{random} & recarga donde está el SKU; si se quedó sin stock, con probabilidad 0,9 se muda a una ubicación libre al azar & un depósito sin plan; así se generaron los datos \\
\code{objective\_greedy} & a la ubicación libre o propia que menos sube $C$ con el $\lambda$ dado & una decisión en línea guiada por el objetivo \\
\bottomrule
\end{tabularx}
\caption{Políticas de reposición.}
\end{table}

\code{objective\_greedy} evalúa, para cada ubicación candidata $\ell$ del SKU $s$,
\[
\text{costo}(\ell) = \lambda\, f_s\, c_\ell + (1-\lambda)\, 2 \sum_{j \in N(s)} a_{sj}\,
\bar d(\ell, j),
\]
con $\bar d(\ell, j)$ la distancia media de $\ell$ a las ubicaciones que ocupa el vecino $j$.

\code{plan\_aisle} tiene dos variantes. Con \code{trigger="always"} mueve al SKU en cualquier
reposición, aunque le queden unidades (las que quedan se pican en la ubicación vieja, que se
libera al vaciarse); con \code{target="location"} apunta a la ubicación exacta del plan, o a la
libre más cercana caminando, en vez de al pasillo. Bajo el recorrido en S la caminata depende
solo de los pasillos, y por eso la variante por defecto apunta al pasillo.

\begin{hiper}{simulación}
\begin{itemize}
  \item Punto de partida (\code{start}): \code{real}, el depósito al corte (por defecto), o
        \code{plan}, el plan aplicado de una vez (el techo).
  \item Política de reposición: \code{plan\_aisle} (por defecto), \code{home}, \code{random} o
        \code{objective\_greedy}.
  \item \code{rule\_trigger} y \code{rule\_target} de \code{plan\_aisle}: \code{stockout} y
        \code{aisle} por defecto.
  \item \code{relocation\_probability} de \code{random}: por defecto 0,9, la de los datos.
  \item \code{lam} de \code{objective\_greedy}: se pasa al simulador; en los experimentos 0,5.
  \item \code{routing}: \code{snake} o \code{s\_shape} (capítulo~\ref{cap:ruteo}).
  \item \code{seed}: la semilla del sorteo de \code{random}.
\end{itemize}
\end{hiper}

\section{El ciclo completo}

\begin{figure}[ht]
\centering
\begin{tikzpicture}[node distance=6mm and 6mm]
  \node[caja, fill=sim!10, draw=sim] (ini) {estado real\\\scriptsize al corte};
  \node[caja, fill=sim!10, draw=sim, right=of ini] (ev) {próximo evento};
  \node[caja, fill=sim!10, draw=sim, above right=4mm and 6mm of ev] (pick) {línea de pick\\\scriptsize consume stock, anota la ubicación};
  \node[caja, fill=plan!12, draw=plan, below right=4mm and 6mm of ev] (rep) {reposición\\\scriptsize la política elige el destino};
  \node[caja, fill=sim!10, draw=sim, right=of pick] (cierre) {fin del batch\\\scriptsize caminata con la política de ruteo};
  \node[caja, fill=sim!25, draw=sim, below=14mm of cierre] (res) {\code{SimulationResult}};
  \draw[flecha] (ini) -- (ev);
  \draw[flecha] (ev) |- (pick);
  \draw[flecha] (ev) |- (rep);
  \draw[flecha] (pick) -- (cierre);
  \draw[flecha] (rep.west) -| ([xshift=-3mm]ev.south);
  \draw[flecha] (cierre) -- node[right, rotulo]{al terminar\\el test} (res);
\end{tikzpicture}
\caption{El ciclo de la simulación.}
\label{fig:ciclo}
\end{figure}

\section{Pool por pool, con el depósito del momento}
\label{sec:rolling}

Partiendo del depósito real, el pasillo de un SKU cambia durante el test, así que los batches
no se pueden armar de antemano con los pasillos del plan. La función \code{simulate\_rolling}
corta las órdenes de test en pools de $25\,W$ órdenes consecutivas y, para cada pool:

\begin{enumerate}
  \item lee el pasillo actual de cada SKU, el de su ubicación con más unidades;
  \item re-arma el pool con \code{AisleCoverBatching} usando esos pasillos
        (capítulo~\ref{cap:batching});
  \item simula sus batches y las reposiciones hasta su última línea, y deja el estado para el
        pool siguiente.
\end{enumerate}

Con $W = 1$ los batches son FIFO y no dependen del estado. Cada paso registra cuánto avanzó el
plan: la fracción de la demanda de train cuyos SKUs ya están en su pasillo del plan, y los
pasillos que toca en promedio una orden con los SKUs donde están.

Las medidas de una corrida son dos: la caminata media por batch en todo el test
(\emph{promedio}), que incluye la transición, y la de los batches de los últimos 14 días
(\emph{final}), que muestra hasta dónde se llega. Ambas se comparan con \emph{hoy}: el depósito
real con reposición \code{random}, como se generaron los datos, y batches FIFO.

\section{Salidas y controles}

El resultado de una corrida tiene:

\begin{itemize}
  \item las métricas de la caminata por batch: total, media, mediana y percentil 95;
  \item un registro por batch (hora, distancia, pasillos visitados, líneas, faltantes);
  \item un registro por reposición (destino, si fue mudanza, si fue a la ubicación del plan);
  \item contadores: mudanzas, líneas y unidades faltantes, y la fracción de SKUs que terminan
        fuera de su ubicación del plan.
\end{itemize}

Controles que se verificaron:

\begin{itemize}
  \item Con la política \code{home} y un plan que entra en el depósito, ningún SKU se mueve y
        la caminata simulada es idéntica a la que calcula el algoritmo de re-batching
        (diferencia 0 en todas las evaluaciones con \code{home} y recorrido en S del notebook 05).
  \item Con el plan aplicado y \code{home}, la simulación pool por pool da exactamente lo mismo
        que armar todos los batches de una vez (notebook \code{correct\_base/01\_transicion}).
  \item \code{reslot\_state} conserva las unidades de cada SKU.
  \item En la v1, reproducir la historia completa da el mismo stock final que un cálculo
        independiente, y los quiebres de stock del test coinciden con los de los datos.
  \item En \code{medium} y \code{medium-alpha-1.3} no hay faltantes. En \code{big} hay entre 1
        y 25 líneas sin stock por corrida, de unas 317.000, y aparecen igual con FIFO: vienen
        del borde entre train y test.
\end{itemize}

% =============================================================================
\chapter{Arquitectura del código}
\label{cap:arquitectura}

El código es un paquete de Python, \code{abs\_affinity\_based\_slotting}, en
\code{src/}. Tiene unas 5.200 líneas repartidas en diez subpaquetes, uno por concepto.

\section{Las capas}

Cada subpaquete depende solo de los de capas anteriores (figura~\ref{fig:capas}).

\begin{figure}[ht]
\centering
\begin{tikzpicture}[node distance=4mm and 4mm, pk/.style={caja, minimum width=26mm, font=\footnotesize}]
  \node[pk, fill=datos!15, draw=datos] (data) {\code{data}\\\scriptsize lectura y corte};
  \node[pk, fill=modelo!15, draw=modelo, above left=7mm and -6mm of data] (demand) {\code{demand}\\\scriptsize $f$, co-ocurrencia, $A$};
  \node[pk, fill=modelo!15, draw=modelo, above right=7mm and -6mm of data] (wh) {\code{warehouse}\\\scriptsize $c$, distancias, $D$};
  \node[pk, fill=modelo!25, draw=modelo, above=19mm of data] (slot) {\code{slotting}\\\scriptsize instancia, plan, objetivo};
  \node[pk, fill=plan!15, draw=plan, above left=7mm and -2mm of slot] (clu) {\code{clustering}\\\scriptsize agrupamientos};
  \node[pk, fill=plan!15, draw=plan, above right=7mm and -2mm of slot] (meth) {\code{methods}\\\scriptsize métodos de slotting};
  \node[pk, fill=batch!15, draw=batch, above=7mm of meth] (bat) {\code{batching}\\\scriptsize FIFO, re-batching};
  \node[pk, fill=sim!15, draw=sim, above=19mm of slot] (eval) {\code{evaluation}\\\scriptsize ruteo, métricas};
  \node[pk, fill=sim!15, draw=sim, above left=7mm and -2mm of eval] (simu) {\code{simulation}\\\scriptsize estado, eventos, políticas};
  \node[pk, fill=exp!25, draw=exp, above=19mm of eval] (exp) {\code{experiments}\\\scriptsize tareas, ejecución};
  \draw[flecha] (data) -- (demand); \draw[flecha] (data) -- (wh);
  \draw[flecha] (demand) -- (slot); \draw[flecha] (wh) -- (slot);
  \draw[flecha] (slot) -- (clu); \draw[flecha] (slot) -- (meth); \draw[flecha] (clu) -- (meth);
  \draw[flecha] (meth) -- (bat);
  \draw[flecha] (slot) -- (eval); \draw[flecha] (eval) -- (simu); \draw[flecha] (slot.north west) to[bend left=20] (simu.south);
  \draw[flecha] (simu) -- (exp); \draw[flecha] (bat) -- (exp); \draw[flecha] (meth.east) to[bend right=30] (exp.east);
\end{tikzpicture}
\caption{Subpaquetes y dependencias. Cada flecha va de un subpaquete a otro que lo usa.}
\label{fig:capas}
\end{figure}

\begin{table}[ht]
\centering\small
\begin{tabularx}{\textwidth}{@{}l L L@{}}
\toprule
subpaquete & recibe & devuelve \\
\midrule
\code{data} & las cinco tablas & tablas en memoria; corte train/test \\
\code{demand} & picking de train & demanda $f$; co-ocurrencia; afinidad $A$ \\
\code{warehouse} & stock, distancias, coordenadas & costo $c$; matrices de distancia $D$ \\
\code{slotting} & lo anterior & \code{SlottingInstance}; \code{Assignment}; objetivo y deltas \\
\code{clustering} & una instancia & un grupo por SKU \\
\code{methods} & una instancia & un plan (\code{Assignment}) \\
\code{batching} & picking de un período y el pasillo de cada SKU & el mismo picking con otros \code{batch\_id} \\
\code{evaluation} & bays a visitar & caminata de un batch; métricas \\
\code{simulation} & estado, eventos, política & \code{SimulationResult} \\
\code{experiments} & tareas & resultados en disco \\
\bottomrule
\end{tabularx}
\caption{Qué entra y qué sale de cada subpaquete.}
\end{table}

\section{Registros}
\label{sec:registros}

Cada familia de piezas intercambiables tiene un \emph{registro}: un diccionario de nombre a
implementación. El código pide las piezas por nombre y no depende de las clases concretas.
Agregar una variante es escribir una clase y registrarla con un decorador.

\begin{lstlisting}
@method_registry.register("aisle_fill")
class AisleFillSlotting: ...

metodo = method_registry.get("swap_search")(lam=0.5, neighborhood="bay")
plan = metodo.solve(instancia)
\end{lstlisting}

\begin{table}[ht]
\centering\small
\begin{tabularx}{\textwidth}{@{}l l L@{}}
\toprule
registro & subpaquete & nombres registrados \\
\midrule
\code{affinity\_registry} & \code{demand} & \code{cooccurrence}, \code{jaccard}, \code{cosine} \\
\code{filter\_registry} & \code{demand} & \code{top\_k}, \code{mutual\_top\_k}, \code{threshold} \\
\code{distance\_registry} & \code{warehouse} & \code{walking}, \code{aisle} \\
\code{clustering\_registry} & \code{clustering} & \code{merchant}, \code{demand\_class} \\
\code{method\_registry} & \code{methods} & \code{current}, \code{random}, \code{demand\_greedy}, \code{linear\_assignment}, \code{aisle\_fill}, \code{merchant\_aisle\_fill}, \code{swap\_search}, \code{bilevel}, \code{exact\_qap} \\
\code{batching\_registry} & \code{batching} & \code{fifo}, \code{aisle\_cover} \\
\code{routing\_registry} & \code{evaluation} & \code{snake}, \code{s\_shape} \\
\code{replenishment\_policy\_registry} & \code{simulation} & \code{home}, \code{random}, \code{objective\_greedy} \\
\bottomrule
\end{tabularx}
\caption{Los registros del sistema.}
\end{table}

\section{El árbol del paquete}

\begin{lstlisting}[language={}, numbers=none, basicstyle=\ttfamily\scriptsize]
src/abs_affinity_based_slotting/
  config.py           dataset_dir(version, size); constantes; licencia del solver
  registry.py         el mecanismo nombre -> implementacion
  data/               loaders.py (las 5 tablas), split.py (corte train/test)
  demand/             sku_demand.py (f), cooccurrence.py (n_ij, s_i),
                      affinity.py (metricas), filter.py (filtros)
  warehouse/          locations.py (universo), distances.py (D: walking, aisle),
                      costs.py (c: distancia al dock)
  slotting/           instance.py (SlottingInstance), build.py (build_full_instance),
                      assignment.py (Assignment), objective.py (C y deltas)
  clustering/         merchant.py, abc.py (clases A/B/C)
  methods/            current, demand_greedy, linear_assignment, aisle_fill,
                      local_search (swap_search), two_stage (bilevel), exact
  batching/           fifo.py, aisle_cover.py (re-batching), waits.py (espera)
  evaluation/         routing.py (snake, s_shape), routes.py, metrics.py, evaluator.py
  simulation/         state.py, events.py, initial_state.py, simulator.py, results.py,
                      home.py, random_storage.py, objective_greedy.py,
                      plan_aisle.py (la regla), rolling.py (pool por pool)
  experiments/        specs.py (tareas), runner.py (ejecucion en paralelo),
                      registry.py (registro de experimentos)
\end{lstlisting}

\section{Convenciones}

\begin{itemize}
  \item \textbf{Núcleo numérico por posiciones.} En \code{slotting}, \code{simulation} y
        \code{batching} se trabaja con índices enteros, numpy y scipy. Pandas aparece solo en
        los bordes: al leer datos y al devolver resultados.
  \item \textbf{Nombres por lo que hace, no por la herramienta.} El método exacto se llama
        \code{exact} y el entorno del solver \code{make\_solver\_env}; si cambia el solver, los
        nombres siguen valiendo.
  \item \textbf{Un subpaquete por concepto.}
  \item \textbf{Falla temprana.} La instancia valida dimensiones, valores y simetría al
        construirse; un dato inconsistente lanza una excepción. La excepción a la regla son los
        faltantes de stock en la simulación, que son dinámica legítima y se cuentan.
  \item \textbf{Los datos y los resultados no se versionan} salvo la v1; los resultados de los
        experimentos van a \code{results/}, fuera de git.
\end{itemize}

\section{Un recorrido de punta a punta}

El fragmento siguiente arma la instancia de \code{medium}, calcula un plan, re-arma los
batches de test con ventana 4 y simula con recorrido en S. Es la misma tubería de los
notebooks 04 y 05.

\begin{lstlisting}
ds = WarehouseDataLoader(dataset_dir("v2", "medium")).load_all()
split = split_picking_events(ds.picking_events, test_size=0.2)
inst = build_full_instance(split.train, ds.initial_stock, ds.distances,
                           affinity_metric="jaccard", cooccurrence_col="order_id",
                           affinity_filter=TopKFilter(k=10), normalize=True,
                           pair_distance_model="aisle", coordinates=ds.coordinates)

# estado real al corte
replen_train, replen_test = split_replenishments(ds.replenishment_events, split.train,
                                                 ds.initial_stock, split.cutoff)
state0 = state_from_initial_stock(inst, ds.initial_stock)
state_cut, _ = state_after_history(state0, build_event_stream(split.train, replen_train, inst), inst)

# plan y re-batching
plan = method_registry.get("merchant_aisle_fill")(ds.coordinates).solve(inst)
sim = Simulator.from_tables(inst, ds.coordinates, routing="s_shape")
loc = np.array([inst.location_index(plan.location_of(k)) for k in inst.sku_ids])
sku_aisle = pd.Series(sim.routing.bay_aisle[inst.location_bay[loc]], index=inst.sku_ids)
batching = AisleCoverBatching(sku_aisle, sim.routing.mask_costs(), window=4)
picking = batching.make_batches(split.test)

# simulacion
res = sim.run(reslot_state(state_cut, plan, inst), build_event_stream(picking, replen_test, inst),
              replenishment_policy_registry.get("home")(), planned=plan)
print(res.route_metrics.mean_batch_distance)
\end{lstlisting}

% =============================================================================
\chapter{Los experimentos}
\label{cap:experimentos}

\section{Los notebooks}

Cada experimento es un notebook numerado en \code{notebooks/data2/}. Contiene la formulación,
el código que lo corre y una lectura de los resultados. Los notebooks de la v1 quedaron en
\code{notebooks/data1/}. Los de la v2 se ordenan por régimen de simulación: los que describen
los datos y el código quedan en \code{data2/}; los que aplican el plan de una vez, que miden el
techo, en \code{old\_base/}; y los del régimen vigente, que parten del depósito real, en
\code{correct\_base/}.

\begin{table}[ht]
\centering\small
\begin{tabularx}{\textwidth}{@{}l L@{}}
\toprule
notebook & contenido \\
\midrule
\code{00\_el\_deposito.md} & el depósito contado en palabras \\
\code{01\_EDA\_v2} & análisis exploratorio de la v2: volumen, demanda, co-ocurrencia \\
\code{02\_el\_codigo} & recorrido por el código: cada clase, sus valores y las piezas intercambiables; el exacto en subinstancias \\
\code{06\_datos\_medium} & los datos de \code{medium}, hecho por hecho, para el capítulo del depósito de la tesis \\
\midrule
\code{old\_base/03\_techo\_pasillos} & recorrido en S contra el evaluador viejo; problema de asignar SKUs a pasillos para minimizar pasillos por batch \\
\code{old\_base/04\_rebatching} & re-batching paso a paso, con la grilla de slotting por ventana \\
\code{old\_base/05\_combinatoria} & 138 tareas y 948 evaluaciones: depósito, slotting, ventana, recorrido, frecuencia, parámetros del QAP y reposición \\
\code{old\_base/06\_experimento\_ganador} & el plan por merchant con re-batching, con todos sus modelos y parámetros \\
\midrule
\code{correct\_base/01\_transicion} & el plan entrando por las mudanzas: semana a semana, contra hoy y contra el techo \\
\code{correct\_base/02\_combinatoria} & la grilla del régimen vigente: slotting, ventana, regla, frecuencia, QAP, batching y recorrido \\
\bottomrule
\end{tabularx}
\caption{Los notebooks de la v2.}
\end{table}

\section{El ejecutor de experimentos}

El subpaquete \code{experiments} corre grillas grandes de configuraciones. Describe cada corrida
con tres especificaciones y una tarea (figura~\ref{fig:tarea}):

\begin{itemize}
  \item \code{InstanceSpec}: cómo se arma la instancia (depósito, frecuencia, afinidad,
        normalización, $D$, fracción de test).
  \item \code{SlottingSpec}: el método y sus parámetros.
  \item \code{EvalSpec}: cómo se mide el plan (ventana, pasadas del re-batching, recorrido,
        punto de partida, reposición y sus variantes). Con
        \code{start="real"} la evaluación usa \code{simulate\_rolling}; con
        \code{start="plan"} arma los batches de una vez con los pasillos del plan.
  \item \code{Task}: una instancia, un plan y una lista de evaluaciones.
\end{itemize}

\begin{figure}[ht]
\centering
\begin{tikzpicture}[node distance=5mm and 7mm]
  \node[caja, fill=exp!20, draw=exp] (t) {\code{Task}};
  \node[caja, fill=modelo!15, draw=modelo, right=of t] (c) {contexto\\\scriptsize instancia, estado al\\\scriptsize corte, simuladores};
  \node[caja, fill=plan!15, draw=plan, right=of c] (p) {plan\\\scriptsize una vez};
  \node[caja, fill=batch!15, draw=batch, right=of p, yshift=8mm] (e1) {evaluación 1\\\scriptsize re-batching y simulación};
  \node[caja, fill=batch!15, draw=batch, right=of p, yshift=-8mm] (e2) {evaluación $n$};
  \node[caja, fill=exp!20, draw=exp, right=of p, xshift=36mm] (disk) {una fila por\\evaluación\\\scriptsize \code{tareas/<clave>.parquet}};
  \draw[flecha] (t) -- (c); \draw[flecha] (c) -- (p); \draw[flecha] (p) -- (e1); \draw[flecha] (p) -- (e2);
  \draw[flecha] (e1) -- (disk); \draw[flecha] (e2) -- (disk);
  \node[rotulo] at ($(e1)!0.5!(e2)$) {$\vdots$};
\end{tikzpicture}
\caption{Una tarea: el plan, que es lo caro, se calcula una vez y se mide de varias formas.}
\label{fig:tarea}
\end{figure}

\paragraph{Cómo corre.} \code{run\_tasks} recibe la lista de tareas y las corre en paralelo, un
depósito por vez. Cada tarea tiene una clave fija que sale de sus especificaciones. Al terminar
guarda su resultado en un archivo con esa clave, que aparece solo cuando está completo. Si la
corrida se corta, al repetirla se saltean las tareas que ya tienen archivo. Dentro de una tarea,
el re-batching de cada ventana se calcula una vez y se reusa para los dos recorridos.

\paragraph{Memoria y licencia.} Un proceso de \code{big} ocupa unos 3 GB, así que \code{big}
corre de a un proceso; los demás depósitos, de a cuatro. Las tareas que usan el solver
(\code{bilevel}) corren de a una, porque la licencia limita las sesiones simultáneas. El avance
queda en \code{progreso.log}.

\paragraph{Salidas.} Una fila por evaluación con todas las especificaciones y las métricas:
caminata media y percentil 95, pasillos por batch, espera media y máxima, la caminata que
calcula el algoritmo de re-batching (FIFO, construcción y final), faltantes, mudanzas, fracción
fuera del plan y tiempos; con \code{start="real"}, además, la caminata de los últimos 14 días,
la demanda en su pasillo del plan al final y los pasillos por orden al final. Aparte, un
archivo con la caminata de cada batch simulado, para analizar por semana, y otro con el registro
de cada paso de \code{simulate\_rolling} (\code{pasos/}).

\paragraph{Registro de experimentos.} Un experimento es una configuración: la instancia, el plan
y cómo se mide. \code{register\_results} lleva cada una a \code{results/registro.db}. El nombre
es el plan; el dataset, el depósito con su partición train/test; los parámetros, todos los
campos de la configuración, con \code{-} donde la especificación los deja vacíos (los de
\code{swap\_search} y \code{bilevel} en los otros métodos), y también en un
\code{configuracion.json}; los resultados son la caminata promedio y
final, su diferencia contra hoy, pasillos por batch, espera y mudanzas; y la descripción dice la
configuración en palabras y en qué bloques aparece. Una configuración que
aparece en varios bloques da el mismo resultado en todos y se registra una vez, con la lista de
sus bloques; las que ya están no se repiten. La grilla del régimen vigente tiene 807
configuraciones. El registro se mira con
\code{mlflow ui --backend-store-uri sqlite:///results/registro.db}: una tabla con una fila por
configuración, que se puede filtrar y ordenar por cualquier parámetro o resultado.

\paragraph{Claves de las tareas viejas.} Los campos que se agregaron después de la grilla 05
(\code{test\_size}, \code{start}, \code{rule\_trigger}, \code{rule\_target}) quedan fuera de la clave cuando tienen el valor que esa grilla usaba
sin decirlo, así su caché sigue sirviendo. El notebook 05 fija esos valores en su celda de
configuración.

\section{El diseño de la grilla del régimen vigente}

La grilla de \code{correct\_base/02\_combinatoria} busca descartar las opciones que no suman y
elegir el valor más útil de cada factor. Usa la misma idea que la del notebook 05: una grilla
completa con slotting por ventana y barridos de a un factor para el resto
(tabla~\ref{tab:bloques02}). Cada evaluación se compara con \emph{hoy} en dos medidas, el
promedio de todo el test y el final (los últimos 14 días).

\begin{table}[ht]
\centering\small
\begin{tabularx}{\textwidth}{@{}l L L@{}}
\toprule
bloque & qué varía & dónde \\
\midrule
A & 11 planes (6 en \code{big}) $\times$ ventanas 1, 2, 4, 8, 16, con la regla y con el techo; hoy con el vigente & 4 depósitos \\
R & regla \code{plan\_aisle} (\code{stockout}/\code{always} $\times$ \code{aisle}/\code{location}) con 4 planes; \code{objective\_greedy} con $\lambda$ = 0, 0,5 y 1; hoy & 4 depósitos \\
F & las 4 frecuencias con 3 planes & 3 depósitos \\
Q & 14 variantes de \code{swap\_search} ($\lambda$, $D$, métrica, unidad y filtro de la afinidad, normalización, vecindad, plan inicial) & 3 depósitos \\
B & pasadas del re-batching: 0 y 20 & 3 depósitos \\
G & recorrido \code{s\_shape} y \code{snake} & 3 depósitos \\
\bottomrule
\end{tabularx}
\caption{Los bloques de la grilla del régimen vigente.}
\label{tab:bloques02}
\end{table}

\paragraph{Configuración base.} Igual que la del notebook 05, con el régimen vigente: test del
último 40\%, depósito real al empezar, regla \code{plan\_aisle} con \code{stockout} y
\code{aisle}.

\section{El diseño del notebook 05}

Todas las combinaciones de todos los factores son decenas de miles de corridas, y
\code{swap\_search} tarda de 5 minutos a varias horas por plan. El diseño tiene una grilla
completa con los factores que más mueven el resultado y barridos de a un factor para el resto
(tabla~\ref{tab:bloques}). Un barrido parte de una configuración base y cambia un solo factor
por vez: muestra el efecto de ese factor con todo lo demás fijo, pero no cómo interactúan dos
factores barridos.

\begin{table}[ht]
\centering\small
\begin{tabularx}{\textwidth}{@{}l L L r@{}}
\toprule
bloque & qué varía & evaluaciones por plan & tareas \\
\midrule
A, principal & 4 depósitos $\times$ 10 planes (\code{swap\_search} y \code{linear\_assignment} no en \code{big}) & 6 ventanas $\times$ 2 recorridos & 38 \\
B, frecuencia & 3 frecuencias más $\times$ 3 planes $\times$ 4 depósitos & 3 ventanas, recorrido en S & 36 \\
C, QAP & 16 variantes de \code{swap\_search} $\times$ 3 depósitos & 3 ventanas $\times$ 2 recorridos & 48 \\
D, reposición & 4 planes $\times$ 4 depósitos & 2 ventanas $\times$ 3 políticas & 16 \\
\midrule
total & & & 138 \\
\bottomrule
\end{tabularx}
\caption{Los bloques del notebook 05: 138 tareas y 948 evaluaciones. La corrida completa tardó
unas 20 horas, casi todo en las variantes lentas de \code{swap\_search}.}
\label{tab:bloques}
\end{table}

\paragraph{Configuración base.} Frecuencia \code{pick\_lines}; afinidad \code{jaccard} por
orden con \code{top\_k:10}; normalización; $D$ = \code{aisle}; reposición \code{home}. Para
\code{swap\_search}: $\lambda = 0{,}5$, vecindad \code{bay}, plan inicial
\code{demand\_greedy}.

\section{Catálogo de hiperparámetros}

\begin{footnotesize}
\begin{longtable}{@{}>{\raggedright\arraybackslash}p{3.4cm} >{\raggedright\arraybackslash}p{2.4cm} >{\raggedright\arraybackslash}p{4.0cm} >{\raggedright\arraybackslash}p{2.0cm} >{\raggedright\arraybackslash}p{2.6cm}@{}}
\toprule
parámetro & pieza & valores & por defecto & en el notebook 05 \\
\midrule
\endhead
\multicolumn{5}{@{}l}{\textbf{Datos}} \\
\code{dataset\_dir} & \code{config} & \code{v1}; \code{v2} con \code{small}, \code{medium}, \code{big} o una variante & --- & los 4 de la v2 \\
\code{test\_size} & corte & fracción en $(0, 1)$ & 0,4 & 0,2 \\
\multicolumn{5}{@{}l}{\textbf{Instancia}} \\
\code{demand\_metric} & $f$ & \code{pick\_lines}, \code{unique\_orders}, \code{unique\_batches}, \code{total\_units} & \code{pick\_lines} & las 4 \\
\code{cooccurrence\_col} & $A$ & \code{order\_id}, \code{batch\_id} & \code{batch\_id} & las 2 \\
\code{affinity\_metric} & $A$ & \code{jaccard}, \code{cosine}, \code{cooccurrence} & \code{jaccard} & las 3 \\
\code{affinity\_filter} & $A$ & \code{TopKFilter(k)}, \code{MutualTopKFilter(k)}, \code{ThresholdFilter(t)}, ninguno & ninguno & top-5, 10, 20; mutuo 10; ninguno \\
\code{normalize} & $f$ y $A$ & sí, no & no & los 2 \\
\code{pair\_distance\_model} & $D$ & \code{walking}, \code{aisle} & \code{walking} & los 2 \\
\multicolumn{5}{@{}l}{\textbf{Slotting}} \\
método & \code{methods} & ver tabla~\ref{tab:metodos} & --- & 10 planes \\
\code{lam} & \code{swap\_search} & $[0, 1]$ & 0,5 & 0; 0,25; 0,5; 0,75 \\
\code{neighborhood} & \code{swap\_search} & \code{bay}, \code{neighbor} & \code{bay} & los 2 \\
\code{seed\_method} & \code{swap\_search} & cualquier método & \code{demand\_}\allowbreak\code{greedy} & greedy, \code{aisle\_}\allowbreak\code{fill}, \code{merchant\_}\allowbreak\code{aisle\_fill} \\
\code{max\_passes} & \code{swap\_search} & entero & 50 & 50 \\
\code{clustering} & \code{bilevel} & \code{merchant}, \code{demand\_class} & --- & los 2 \\
\code{thresholds} & \code{demand\_class} & cortes de demanda acumulada & $(0{,}8;\ 0{,}95)$ & por defecto \\
\code{zone\_cost} & \code{bilevel} & \code{dock}, \code{aisle\_order} & \code{dock} & los 2 \\
\code{zone\_solver} & \code{bilevel} & cualquier método & --- & \code{demand\_greedy} \\
\code{lam}, \code{time\_limit}, \code{mip\_gap} & \code{exact\_qap} & --- & 0,5; sin límite & no se usa \\
\multicolumn{5}{@{}l}{\textbf{Batching}} \\
\code{batch\_size} & \code{batching} & entero & 25 & 25 \\
\code{window} & \code{aisle\_cover} & entero $\ge 1$ & 1 & 1, 2, 4, 8, 16, 32 \\
\code{max\_passes} & \code{aisle\_cover} & entero $\ge 0$ & 20 & 20 \\
\multicolumn{5}{@{}l}{\textbf{Simulación}} \\
\code{routing} & simulador & \code{snake}, \code{s\_shape} & \code{snake} & los 2 \\
\code{start} & \code{EvalSpec} & \code{real}, \code{plan} & \code{real} & \code{plan} \\
política de reposición & simulador & \code{plan\_aisle}, \code{home}, \code{random}, \code{objective\_greedy} & \code{plan\_aisle} & las 3 últimas \\
\code{trigger} & \code{plan\_aisle} & \code{stockout}, \code{always} & \code{stockout} & no existía \\
\code{target} & \code{plan\_aisle} & \code{aisle}, \code{location} & \code{aisle} & no existía \\
\code{relocation\_}\allowbreak\code{probability} & \code{random} & $[0, 1]$ & 0,9 & 0,9 \\
\code{lam} & \code{objective\_}\allowbreak\code{greedy} & $[0, 1]$ & 0,5 & 0,5 \\
\code{seed} & simulador & entero & 0 & 0 \\
\bottomrule
\caption{Todos los hiperparámetros del sistema.}
\label{tab:catalogo}
\end{longtable}
\end{footnotesize}

\section{Cómo reproducir}

\begin{lstlisting}[language=bash, numbers=none]
# entorno
source .venv/bin/activate

# un notebook, ejecutado en el lugar; las grillas retoman desde su carpeta de results/
jupyter nbconvert --to notebook --execute --inplace notebooks/data2/old_base/05_combinatoria.ipynb
jupyter nbconvert --to notebook --execute --inplace --ExecutePreprocessor.timeout=-1 \
    notebooks/data2/correct_base/02_combinatoria.ipynb

# este documento
make manual
\end{lstlisting}

La licencia del solver va en el archivo \code{.env} de la raíz, con tres variables:
\begin{itemize}
  \item \code{SOLVER\_LICENSE\_ACCESS\_ID},
  \item \code{SOLVER\_LICENSE\_SECRET},
  \item \code{SOLVER\_LICENSE\_ID}.
\end{itemize} Solo \code{bilevel} y \code{exact\_qap} usan el solver; el
resto del sistema no lo necesita.

% =============================================================================
\chapter{Resultados}
\label{cap:resultados}

Las cifras de este capítulo salen de tres notebooks de \code{correct\_base/}: la grilla
\code{02\_combinatoria}, \code{04\_experimento\_ganador} y \code{05\_convergencia}. Todas son la caminata media por
batch con recorrido en S, simulando el período de test (72 días, unas 10 semanas) desde el
depósito real al corte, como diferencia porcentual contra \emph{hoy}: el plan vigente, reposición \code{random} y batches
FIFO, en el mismo depósito. El \emph{final} es la caminata de los últimos 14 días y muestra
hasta dónde llega cada configuración; el \emph{promedio}, la de todo el test, incluye la
transición. El \emph{techo} es el plan aplicado de una vez con reposición \code{home}: las
cifras de los notebooks de \code{old\_base/} son techos.

\section{Lo que dicen los datos}

La co-ocurrencia por orden no se distingue de la que darían órdenes armadas al azar según la
popularidad (sección~\ref{sec:hechos}). No hay afinidad plantada entre SKUs; la única
estructura conjunta es que una orden nunca mezcla merchants.

\section{El experimento ganador}

El ganador combina el plan por merchant (\code{merchant\_aisle\_fill}), la regla de reposición
\code{plan\_aisle} con sus valores base (mover solo al agotarse, a un hueco del pasillo del
plan) y re-batching con ventana 4 y 20 pasadas. El plan entra solo por las mudanzas de las
reposiciones. En \code{medium} gana 27,4\% al final y 21,5\% en promedio, algo más de la mitad
de su techo, 46,3\% (tabla~\ref{tab:ganador}). Al final del test el 88\% de la demanda está en
su pasillo del plan y una orden toca 1,34 pasillos, contra 1,03 con el plan completo: lo que
falta son los SKUs que casi no se agotan y por eso casi no se mudan. En los 72 días del test
la caminata no se estabiliza: entre las semanas 5 y 9 el ganador sigue bajando unas 170
pulgadas por semana, el doble de lo que explica la variación semanal del techo, y al final
recorrió la mitad del camino entre un plan aleatorio y el techo (notebook 05).

\begin{table}[ht]
\centering\small
\begin{tabular}{@{}lrrrr@{}}
\toprule
corrida & promedio & final & mudanzas & espera (min) \\
\midrule
hoy & 0,0 & 0,0 & 13.608 & 17 \\
hoy con re-batching & $-4{,}1$ & $-4{,}2$ & 12.570 & 113 \\
solo la regla & $-0{,}2$ & $-0{,}4$ & 3.366 & 17 \\
aleatorio & $-8{,}1$ & $-9{,}3$ & 3.328 & 113 \\
greedy & $-13{,}7$ & $-17{,}1$ & 4.398 & 113 \\
\textbf{ganador} & $\mathbf{-21{,}5}$ & $\mathbf{-27{,}4}$ & 3.362 & 113 \\
techo & $-47{,}1$ & $-46{,}3$ & 0 & 113 \\
\bottomrule
\end{tabular}
\caption{El ganador y sus piezas en \code{medium}: diferencia porcentual contra hoy, mudanzas
por reposición en todo el test y espera media de una orden. El aleatorio es un plan sorteado
(\code{random}) y el greedy, \code{demand\_greedy}, los dos con la misma regla y el mismo
re-batching que el ganador.}
\label{tab:ganador}
\end{table}

\begin{itemize}
  \item \textbf{El plan vale 20\% contra un plan sin datos.} Simulado igual que un plan
        aleatorio, el ganador camina 19,9\% menos al final y 14,6\% menos en promedio (24,3\% y
        18,0\% en \code{big}); tres sorteos del aleatorio difieren en menos de un punto. Contra el
        greedy simulado igual gana 12,4\% al final (6,7\% en \code{big}): el greedy pone la mitad
        de la demanda en el pasillo 1, y la regla no logra vaciarlo de los SKUs que no van ahí. Contra
        hoy, el aleatorio ya gana 9,3\%: 4,2 puntos del re-batching y 5,1 de la regla, que muda
        menos SKUs, así que los batches se arman con pasillos que siguen siendo ciertos mientras
        se recorren.
  \item \textbf{Hacen falta las dos piezas.} Sin re-armar los batches, el plan gana 0,4\%; sin
        plan, re-armar los batches gana 4,2\%.
  \item \textbf{La regla no agrega mudanzas:} hace 3.362 en el test contra 13.608 de hoy, porque
        mueve un SKU solo cuando se agota y no está en su pasillo, cuando hoy se muda casi
        siempre que se agota.
  \item \textbf{El costo es la espera:} 113 minutos de media con ventana 4, contra 17 con FIFO.
        Con el plan por merchant, la ventana 2 da $-13{,}8\%$ con 49 minutos y la 16,
        $-37{,}1\%$ con 496.
\end{itemize}

En \code{medium-alpha-1.3} y en \code{big} el patrón se repite: el ganador llega a algo más de
la mitad de su techo (tabla~\ref{tab:depositos}). En \code{small} ningún plan gana: con 4
pasillos, mientras quede demanda fuera de su pasillo casi cualquier batch entra a los 4, y
además sale un batch cada 240 minutos, así que con ventana 4 una orden espera 843 minutos de
media.

\begin{table}[ht]
\centering\small
\begin{tabular}{@{}lrrrr@{}}
\toprule
depósito & promedio & final & techo (final) & espera con $W = 4$ (min) \\
\midrule
\code{small} & $-0{,}1$ & $-0{,}2$ & $-28{,}3$ & 843 \\
\code{medium} & $-21{,}5$ & $-27{,}4$ & $-46{,}3$ & 113 \\
\code{medium-alpha-1.3} & $-23{,}3$ & $-28{,}9$ & $-48{,}1$ & 113 \\
\code{big} & $-30{,}8$ & $-36{,}5$ & $-59{,}8$ & 51 \\
\bottomrule
\end{tabular}
\caption{El ganador contra hoy en los cuatro depósitos, y su techo.}
\label{tab:depositos}
\end{table}

\section{Slotting}

Con FIFO casi ningún plan gana: en \code{medium} ninguno pasa de 1\%, porque un batch FIFO
junta 25 órdenes de casi todos los pasillos y entra a los 8 igual. En \code{big}, con 16
pasillos, los planes por demanda llegan a 8,0\%. Con re-batching ganan los dos planes que
agrupan por merchant, en todos los depósitos y desde la ventana 2 (tabla~\ref{tab:principal} y
figura~\ref{fig:principal}). Los planes por demanda y los otros tres binivel quedan entre
$-16{,}8\%$ y $-17{,}4\%$ en \code{medium} con ventana 4, y \code{swap\_search} en el medio. El
vigente con la regla ya gana 11,0\%: dejar de mudar SKUs al azar vale algo por sí solo.

\begin{table}[ht]
\centering\small
\begin{tabular}{@{}llrrrrr@{}}
\toprule
 & & \multicolumn{5}{c}{ventana $W$} \\
\cmidrule(l){3-7}
depósito & plan & 1 & 2 & 4 & 8 & 16 \\
\midrule
\code{medium} & hoy con re-batching & 0,0 & $-0{,}5$ & $-4{,}2$ & $-5{,}2$ & $-5{,}0$ \\
 & vigente & 0,0 & $-1{,}7$ & $-11{,}0$ & $-15{,}1$ & $-18{,}0$ \\
 & \code{demand\_greedy} & $-0{,}6$ & $-11{,}1$ & $-17{,}1$ & $-21{,}7$ & $-24{,}0$ \\
 & \code{aisle\_fill} & $-1{,}0$ & $-10{,}5$ & $-16{,}8$ & $-19{,}9$ & $-22{,}6$ \\
 & \code{swap\_search} & $-0{,}3$ & $-12{,}5$ & $-21{,}7$ & $-27{,}1$ & $-30{,}4$ \\
 & \code{bilevel merchant/aisle\_order} & $-0{,}3$ & $-13{,}2$ & $-26{,}6$ & $-33{,}2$ & $-36{,}8$ \\
 & \code{merchant\_aisle\_fill} & $-0{,}4$ & $-13{,}8$ & $\mathbf{-27{,}4}$ & $-34{,}1$ & $-37{,}1$ \\
 & techo de \code{merchant\_aisle\_fill} & $-3{,}0$ & $-32{,}9$ & $-46{,}3$ & $-57{,}9$ & $-61{,}9$ \\
\midrule
\code{big} & hoy con re-batching & 0,0 & & $-13{,}4$ & & $-23{,}8$ \\
 & vigente & 0,0 & & $-16{,}6$ & & $-31{,}8$ \\
 & \code{demand\_greedy} & $-8{,}0$ & & $-32{,}0$ & & $-40{,}8$ \\
 & \code{aisle\_fill} & $-7{,}4$ & & $-29{,}1$ & & $-38{,}4$ \\
 & \code{bilevel merchant/aisle\_order} & $-5{,}7$ & & $-36{,}5$ & & $-53{,}0$ \\
 & \code{merchant\_aisle\_fill} & $-6{,}0$ & & $-36{,}5$ & & $-53{,}1$ \\
 & techo de \code{merchant\_aisle\_fill} & $-18{,}9$ & & $-59{,}8$ & & $-75{,}8$ \\
\bottomrule
\end{tabular}
\caption{Final contra hoy con la regla base. En \code{big} se corrieron las ventanas 1, 4 y 16. La espera media con cada ventana está en la
tabla~\ref{tab:espera}.}
\label{tab:principal}
\end{table}

\begin{figure}[ht]
\centering
\begin{tikzpicture}
\begin{axis}[width=8.6cm, height=7cm, xmode=log, log basis x=2, grid=major, grid style={gray!20}, xlabel={espera media (horas)}, ylabel={final contra hoy (\%)}, tick label style={font=\scriptsize}, label style={font=\small}, ymin=-68, ymax=4, xtick={0.25,0.5,1,2,4,8}, xticklabels={{0,25},{0,5},1,2,4,8},
  legend pos=outer north east, legend style={font=\scriptsize, draw=none}, legend cell align=left]
  \addplot[gris, dashed, mark=o, mark size=1.2pt] coordinates {(0.283,0) (0.817,-0.5) (1.883,-4.2) (4.017,-5.2) (8.267,-5.0)};
  \addlegendentry{hoy con re-batching}
  \addplot[gris, mark=*, mark size=1.2pt] coordinates {(0.283,0) (0.817,-1.7) (1.883,-11.0) (4.017,-15.1) (8.267,-18.0)};
  \addlegendentry{vigente}
  \addplot[datos, mark=*, mark size=1.2pt] coordinates {(0.283,-0.6) (0.817,-11.1) (1.883,-17.1) (4.017,-21.7) (8.267,-24.0)};
  \addlegendentry{\code{demand\_greedy}}
  \addplot[plan, mark=*, mark size=1.2pt] coordinates {(0.283,-1.0) (0.817,-10.5) (1.883,-16.8) (4.017,-19.9) (8.267,-22.6)};
  \addlegendentry{\code{aisle\_fill}}
  \addplot[modelo, mark=triangle*, mark size=1.4pt] coordinates {(0.283,-0.3) (0.817,-12.5) (1.883,-21.7) (4.017,-27.1) (8.267,-30.4)};
  \addlegendentry{\code{swap\_search}}
  \addplot[batch, dashed, mark=square*, mark size=1.2pt] coordinates {(0.283,-0.3) (0.817,-13.2) (1.883,-26.6) (4.017,-33.2) (8.267,-36.8)};
  \addlegendentry{\code{bilevel merchant/aisle\_order}}
  \addplot[sim, mark=*, mark size=1.4pt] coordinates {(0.283,-0.4) (0.817,-13.8) (1.883,-27.4) (4.017,-34.1) (8.267,-37.1)};
  \addlegendentry{\code{merchant\_aisle\_fill}}
  \addplot[sim, dotted, thick, mark=x, mark size=2pt] coordinates {(0.283,-3.0) (0.817,-32.9) (1.883,-46.3) (4.017,-57.9) (8.267,-61.9)};
  \addlegendentry{techo de \code{merchant\_aisle\_fill}}
\end{axis}
\end{tikzpicture}
\caption{\code{medium}: final contra hoy y espera media de las órdenes. Cada punto es una
ventana: 1, 2, 4, 8 y 16 batches. El techo del plan por merchant queda entre 19 y 25 puntos
debajo de lo que se logra partiendo del depósito real.}
\label{fig:principal}
\end{figure}

En \code{small} ningún plan gana más de 0,3\%, aunque el techo del plan por merchant con
ventana 4 es 28,3\%.

\section{La regla de reposición}

La regla base es la mejor o empata con todos los planes, y es la que menos mueve
(tabla~\ref{tab:regla}). Mover también cuando al SKU le quedan unidades (\code{always}) o
apuntar a la ubicación exacta del plan (\code{location}) empata con el plan por merchant y
pierde 3 a 6 puntos con los planes por demanda, con más mudanzas; las dos cosas juntas casi
borran la ganancia y mudan el doble que hoy. En \code{big}, \code{always/aisle} gana 0,8 puntos
más que la base con 58\% más mudanzas.

\begin{table}[ht]
\centering\small
\begin{tabular}{@{}lrrrr@{}}
\toprule
regla & \code{demand\_greedy} & \code{aisle\_fill} & \code{merchant\_aisle\_fill} & mudanzas \\
\midrule
\code{stockout/aisle} (base) & $-17{,}1$ & $-16{,}8$ & $-27{,}4$ & 3.362 \\
\code{stockout/location} & $-11{,}1$ & $-12{,}5$ & $-27{,}1$ & 12.118 \\
\code{always/aisle} & $-14{,}1$ & $-12{,}4$ & $-27{,}1$ & 5.172 \\
\code{always/location} & $-2{,}2$ & $-2{,}5$ & $-3{,}5$ & 26.424 \\
\midrule
\multicolumn{5}{@{}l}{\emph{sin plan}} \\
\code{objective\_greedy}, $\lambda = 0$ & \multicolumn{3}{c}{$-20{,}2$} & 8.654 \\
\code{objective\_greedy}, $\lambda = 0{,}5$ & \multicolumn{3}{c}{$-15{,}6$} & 5.545 \\
\code{objective\_greedy}, $\lambda = 1$ & \multicolumn{3}{c}{$-11{,}7$} & 1.809 \\
vigente con la regla & \multicolumn{3}{c}{$-11{,}0$} & 0 \\
\code{random} (hoy con re-batching) & \multicolumn{3}{c}{$-4{,}2$} & 12.570 \\
\bottomrule
\end{tabular}
\caption{Final contra hoy en \code{medium} con ventana 4. Las mudanzas
son las de \code{merchant\_aisle\_fill}, o las de la política sin plan.}
\label{tab:regla}
\end{table}

Sin plan, la mejor reposición es \code{objective\_greedy} con la afinidad: con $\lambda = 0$,
que decide solo por afinidad, gana 20,2\% en \code{medium}, y con $\lambda = 1$, solo por
demanda, 11,7\%. La afinidad de los datos, que es el merchant, sirve aun sin plan, pero queda
por debajo del plan por merchant con la regla.

\section{Los parámetros del QAP}

Ninguna variante de \code{swap\_search} supera a \code{merchant\_aisle\_fill}
(tabla~\ref{tab:qap}). La mejor arranca del plan por merchant, y los intercambios que bajan el
objetivo del QAP sacan algunas órdenes de su pasillo (1,08 pasillos por orden contra 1,03).
Arrancando de \code{demand\_greedy}, lo que más suma es lo que acerca el plan al de merchant:
20 vecinos por SKU o \code{cosine} dejan las órdenes en 1,21 a 1,27 pasillos, contra 1,53 de la
base. $\lambda$ entre 0 y 0,75 da casi lo mismo. Empeoran la vecindad \code{neighbor}, la
co-ocurrencia por batch y no normalizar. \code{medium-alpha-1.3} da un orden parecido, y en
\code{small} ninguna variante gana más de 0,5\% partiendo del depósito real.

\begin{table}[ht]
\centering\small
\begin{tabular}{@{}lrrr@{}}
\toprule
variante (\code{medium}) & $W = 4$ & $W = 16$ & minutos \\
\midrule
base ($\lambda = 0{,}5$, \code{bay}, \code{jaccard}, \code{top\_k:10}, $D$ = \code{aisle}) & $-21{,}7$ & $-30{,}4$ & 280 \\
plan inicial \code{merchant\_aisle\_fill} & $-25{,}8$ & $-36{,}6$ & 17 \\
plan inicial \code{aisle\_fill} & $-21{,}5$ & $-29{,}6$ & 44 \\
filtro \code{top\_k:20} & $-24{,}2$ & $-34{,}5$ & 103 \\
filtro \code{top\_k:5} & $-19{,}7$ & $-27{,}8$ & 10 \\
filtro \code{mutual\_top\_k:10} & $-21{,}1$ & $-30{,}4$ & 8 \\
métrica \code{cosine} & $-24{,}0$ & $-34{,}2$ & 51 \\
métrica \code{cooccurrence} & $-19{,}7$ & $-25{,}7$ & 223 \\
$\lambda = 0$ & $-22{,}2$ & $-31{,}6$ & 266 \\
$\lambda = 0{,}25$ & $-22{,}3$ & $-31{,}0$ & 34 \\
$\lambda = 0{,}75$ & $-22{,}1$ & $-30{,}0$ & 41 \\
co-ocurrencia por batch & $-17{,}1$ & $-23{,}6$ & 69 \\
sin normalizar & $-17{,}8$ & $-25{,}1$ & 35 \\
vecindad \code{neighbor} & $-17{,}6$ & $-23{,}7$ & 1 \\
$D$ = \code{walking} & $-19{,}7$ & $-28{,}7$ & 24 \\
\midrule
referencia: \code{merchant\_aisle\_fill} solo & $-27{,}4$ & $-37{,}1$ & menos de 1 \\
\bottomrule
\end{tabular}
\caption{Variantes de \code{swap\_search} en \code{medium}: final contra hoy con la regla base,
y tiempo de cálculo del plan (con otro proceso en paralelo).}
\label{tab:qap}
\end{table}

\section{Frecuencia, batching y recorrido}

\begin{itemize}
  \item \textbf{Frecuencia.} Las cuatro definiciones de $f$ quedan a menos de un punto entre sí
        en todos los depósitos y ventanas; se queda \code{pick\_lines}.
  \item \textbf{Batching.} Intercambiar órdenes después de construir los batches suma entre 0,5
        y 2,3 puntos en \code{medium} y \code{medium-alpha-1.3}; se quedan las 20 pasadas.
  \item \textbf{Recorrido.} Con \code{snake}, contra su propio hoy, los planes por demanda ganan
        10\% con FIFO en \code{medium}, porque ponen los SKUs más pedidos al frente de cada
        pasillo. Con ventana 4 el plan por merchant sigue primero, con menos ventaja: $-23{,}2\%$
        contra $-19{,}0\%$ a $-19{,}8\%$.
\end{itemize}

\section{Un techo para la asignación a pasillos}

El notebook \code{old\_base/03\_techo\_pasillos}, del régimen anterior, plantea el problema de
asignar SKUs a pasillos para minimizar los pasillos que visitan los batches, con variables
$x_{ia}$ (el SKU $i$ va al pasillo $a$) y $y_{ba}$ (el batch $b$ entra al pasillo $a$):
\[
\min \sum_{b, a} w_a\, y_{ba} \quad \text{s.a.} \quad \sum_a x_{ia} = 1,\quad
\sum_i x_{ia} \le \kappa_a,\quad |I_b|\, y_{ba} \ge \sum_{i \in I_b} x_{ia},\quad
x, y \in \{0,1\},
\]
con $w_a$ el costo de entrar al pasillo $a$, $\kappa_a$ su capacidad e $I_b$ los SKUs del
batch $b$. Resuelto sobre los batches de test (un \emph{oráculo}, que conoce el futuro), en
\code{big} baja a 7,57 pasillos por batch y $-36{,}1\%$ contra el greedy, con batches FIFO y el
plan aplicado de una vez. Resuelto sobre train, $-8{,}6\%$. La cota del solver queda lejos
(brecha de 90\% o más), así que esos números no son el óptimo. El re-batching resultó una vía
más directa para bajar los pasillos por batch.

\section{Lo que no se mide}

\begin{itemize}
  \item La congestión de pickers en un mismo pasillo: el simulador no la modela.
  \item La espera con horas de llegada reales, que los datos no traen. La ventana se cuenta en
        batches: en \code{small}, ventana 4 son 14 horas de espera.
  \item Merchants con demandas muy distintas: en los datos piden parecido.
  \item Las interacciones entre los factores que se barrieron de a uno.
\end{itemize}

% =============================================================================
\appendix

\chapter{Fragmentos de código}
\label{ap:codigo}

Los fragmentos son las líneas centrales de cada pieza, recortadas donde hay validaciones o
registro de resultados (\code{...}). La versión completa está en \code{src/}.

\section{El objetivo y el cambio por mudanza}
\code{slotting/objective.py}
\begin{lstlisting}
def slotting_cost(assignment, instance, *, lam):
    location_idx = _assignment_location_indices(assignment, instance)
    linear = float(np.dot(instance.demand, instance.location_cost[location_idx]))
    quadratic = _quadratic_cost_from_location_indices(instance, location_idx)
    return lam * linear + (1.0 - lam) * quadratic

def move_cost_delta(assignment, instance, sku, location_id, *, lam):
    ...
    linear_delta = instance.demand[sku_idx] * (
        instance.location_cost[new_idx] - instance.location_cost[old_idx])
    quadratic_delta = 0.0
    for neighbor_idx, weight in row.items():          # vecinos de afinidad de s
        neighbor_loc = ...                            # ubicacion actual del vecino
        quadratic_delta += weight * (
            instance.distance_between_location_indices(new_idx, neighbor_loc)
            - instance.distance_between_location_indices(old_idx, neighbor_loc))
    return lam * linear_delta + (1.0 - lam) * 2.0 * quadratic_delta
\end{lstlisting}

\section{El orden de llenado por pasillos}
\code{methods/aisle\_fill.py}
\begin{lstlisting}
def aisle_fill_location_order(instance, coordinates):
    aisle_of_bay = coordinates.set_index("bay_id")["aisle"].reindex(instance.bay_ids).to_numpy(dtype=float)
    location_aisle = aisle_of_bay[instance.location_bay]
    front_cost = pd.Series(instance.location_cost).groupby(location_aisle).min()   # rango del pasillo
    aisle_rank = front_cost.rank(method="first").reindex(location_aisle).to_numpy()
    return np.lexsort((np.arange(instance.n_locations), instance.location_cost, aisle_rank))

class MerchantAisleFillSlotting:
    def solve(self, instance):
        merchant = pd.Series(instance.merchant_ids)
        merchant_demand = pd.Series(instance.demand).groupby(merchant).sum()
        merchant_rank = merchant.map(merchant_demand.rank(ascending=False, method="first")).to_numpy()
        sku_order = np.lexsort((np.arange(instance.n_skus), -instance.demand, merchant_rank))
        return _assign_in_order(instance, sku_order, aisle_fill_location_order(instance, self._coordinates))
\end{lstlisting}

\section{El recorrido en S}
\code{evaluation/routing.py}
\begin{lstlisting}
def distance_of_aisles(self, aisles):
    visited = np.unique(np.asarray(aisles, dtype=np.int64))
    m = len(visited)
    if m == 0:
        return 0.0
    lateral = float(np.abs(np.diff(self.aisle_x[visited])).sum())
    back_to_front = self.aisle_length if m % 2 == 1 else 0.0
    return (self._to_front(visited[0]) + m * self.aisle_length + lateral
            + back_to_front + self._to_front(visited[-1]))

def mask_costs(self):
    """R_S de todos los conjuntos de pasillos, indexado por mascara de bits."""
    n = self.n_aisles
    masks = np.arange(1 << n, dtype=np.int64)
    bits = ((masks[:, None] >> np.arange(n)) & 1).astype(bool)
    m = bits.sum(axis=1)
    first = np.where(m > 0, bits.argmax(axis=1), 0)
    last = np.where(m > 0, n - 1 - bits[:, ::-1].argmax(axis=1), 0)
    to_front = np.abs(self.aisle_x - self.dock_xy[0]) + (self.front_y - self.dock_xy[1])
    lateral = self.aisle_x[last] - self.aisle_x[first]
    cost = to_front[first] + m * self.aisle_length + lateral \
           + np.where(m % 2 == 1, self.aisle_length, 0.0) + to_front[last]
    cost[0] = 0.0
    return cost
\end{lstlisting}

\section{El pasillo de cada SKU según un plan}
\begin{lstlisting}
loc = np.array([inst.location_index(plan.location_of(k)) for k in inst.sku_ids])
sku_aisle = pd.Series(routing.bay_aisle[inst.location_bay[loc]], index=inst.sku_ids)
\end{lstlisting}

\section{El re-batching}
\code{batching/aisle\_cover.py}. Construcción:
\begin{lstlisting}
def _construct(self, masks):
    n, q = len(masks), self.batch_size
    batch = np.full(n, -1, dtype=np.int64)
    own_cost = self.cost[masks]                       # R_S de cada orden sola
    for k in range(math.ceil(n / q)):
        free = np.flatnonzero(batch < 0)
        seed = free[np.argmax(own_cost[free])]        # semilla: la orden mas cara
        batch[seed], union, size = k, masks[seed], 1
        while size < q:
            free = np.flatnonzero(batch < 0)
            if free.size == 0:
                break
            increase = self.cost[union | masks[free]] - self.cost[union]   # Delta(o | M)
            pick = free[np.argmin(increase)]
            batch[pick], union, size = k, union | masks[pick], size + 1
    return batch
\end{lstlisting}
Intercambios: para cada orden $o$ se evalúa $\delta(o, o')$ contra todas las órdenes a la vez.
\begin{lstlisting}
for o in range(n):
    k1 = batch[o]
    # o va al batch de o' y o' al de o, para todas las o' del pool
    delta = (self.cost[without[o] | masks] + self.cost[without | masks[o]]
             - self.cost[union[k1]] - self.cost[union[batch]])
    delta[batch == k1] = np.inf                       # solo ordenes de otros batches
    j = int(np.argmin(delta))
    if delta[j] >= -1e-9:
        continue
    ...                                               # aplicar el intercambio y actualizar
                                                      # conteos c_ka, uniones y "without"
\end{lstlisting}
\code{without[o]} es la máscara de $M_k^{-o}$, calculada con los conteos por pasillo:
\begin{lstlisting}
without = (((count[batch] - bits) > 0) * pow2).sum(axis=1)
\end{lstlisting}

\section{La espera de las órdenes}
\code{batching/waits.py}
\begin{lstlisting}
def order_waits(picking, window, batch_size=25):
    order = orders_by_arrival(picking)
    position = pd.Series(np.arange(len(order)), index=order)                    # r
    batch = picking.drop_duplicates("order_id").set_index("order_id")["batch_id"].reindex(order)
    pool = position // (window * batch_size)                                      # p
    complete = (position.groupby(pool).transform("max") + 1) / batch_size       # se completa el pool
    oldest = position.groupby(batch).min()
    turn = oldest.groupby(pool.groupby(batch).first()).rank(method="first") - 1  # t
    return complete + batch.map(turn) - position / batch_size
\end{lstlisting}

\section{El ciclo del simulador}
\code{simulation/simulator.py}
\begin{lstlisting}
for batch in batches:
    visited = set()
    for line_ts, sku_idx, qty in zip(batch.line_timestamps, batch.sku_indices, batch.quantities):
        while replen_ptr < len(replens) and replens[replen_ptr].timestamp <= line_ts:
            apply_replen(replens[replen_ptr])          # la politica elige el destino
            replen_ptr += 1
        stops, missing = pick_units(state, int(sku_idx), int(qty))
        visited.update(stops)
        ...                                            # faltantes y "fuera del plan"
    bays = np.unique(instance.location_bay[np.fromiter(visited, dtype=np.int64)])
    distance = self._routing.distance(bays)            # s_shape o snake
    ...
\end{lstlisting}

\section{Una tarea del ejecutor}
\code{experiments/runner.py}
\begin{lstlisting}
def run_task(task, out_dir):
    target = out_dir / "tareas" / f"{task.key}.parquet"
    if target.exists():
        return {"key": task.key, "status": "cached"}   # ya estaba: se saltea
    ctx = context(task.instance)                       # instancia y estado al corte (en cache)
    plan = build_plan(ctx, task.slotting)              # lo caro: una sola vez
    aisle = sku_aisles(ctx, plan)
    rows, batchings = [], {}
    for ev in task.evals:                              # ventana, recorrido, reposicion
        row, log = evaluate(ctx, plan, aisle, ev, batchings)
        rows.append({...})
    ...                                                # escribe a un .tmp y lo renombra
\end{lstlisting}

% =============================================================================
\chapter{Glosario}

\begin{description}[style=nextline, leftmargin=2.6cm, font=\normalfont\bfseries]
  \item[Batch] viaje de picking: 25 órdenes que un picker recoge en una pasada.
  \item[Bay] columna de estantería a un lado de un pasillo; la unidad de distancia.
  \item[Dock] mesa de empaque; inicio y fin de cada recorrido.
  \item[FIFO] armar los batches con órdenes consecutivas, en orden de llegada.
  \item[Franja] tiempo entre un batch y el siguiente.
  \item[Instancia] los datos del problema de slotting: $f$, $c$, $A$ y $D$ con sus índices.
  \item[Merchant] vendedor dueño de un SKU; cada orden es de un solo merchant.
  \item[Mudanza] reposición que cambia la ubicación del SKU.
  \item[Oráculo] un problema resuelto con datos de test; da una cota, no un método.
  \item[Orden] pedido de un cliente: de 2 a 5 SKUs de un mismo merchant.
  \item[Plan] la ubicación de cada SKU que propone un método (\code{Assignment}).
  \item[Pool] grupo de $W \cdot 25$ órdenes consecutivas que el re-batching reparte en $W$
        batches.
  \item[QAP] \emph{Quadratic Assignment Problem}; NP-difícil.
  \item[Re-batching] armar los batches eligiendo qué órdenes van juntas.
  \item[Recorrido en S] el picker recorre entero cada pasillo con picks, alternando el sentido.
  \item[SKU] producto distinto.
  \item[Slotting] asignación de una ubicación a cada SKU.
  \item[Sustituto] función que guía la búsqueda de un método, distinta de la medida con que se
        juzga el resultado.
  \item[Ubicación] hueco físico: bay, estante y bin.
  \item[Ventana] cantidad $W$ de batches FIFO que el re-batching re-arma juntos.
\end{description}

% =============================================================================
\chapter{Referencias}

\begin{itemize}[leftmargin=*, label={}]
  \item Bartholdi, J. J. y Hackman, S. T. (2014). \emph{Warehouse \& Distribution Science}.
        Georgia Institute of Technology.
  \item Deng, Y., Jiang, W., Wang, Y. y Xu, B. (2025). Optimizing order batching and picking
        problems considering the correlation between products under the scattered storage
        mode. \emph{Sustainability}, 17, 1646.
  \item Kofler, M., Beham, A., Wagner, S. y Affenzeller, M. (2014). Affinity based slotting in
        warehouses with dynamic order patterns. En \emph{Advanced Methods and Applications in
        Computational Intelligence}, Springer.
  \item Koopmans, T. C. y Beckmann, M. (1957). Assignment problems and the location of
        economic activities. \emph{Econometrica}, 25(1), 53--76.
  \item Li, J., Moghaddam, M. y Nof, S. Y. (2015). Dynamic storage assignment with product
        affinity and ABC classification: a case study. \emph{International Journal of Advanced
        Manufacturing Technology}.
  \item Loiola, E. M., de Abreu, N. M. M., Boaventura-Netto, P. O., Hahn, P. y Querido, T.
        (2007). A survey for the quadratic assignment problem. \emph{European Journal of
        Operational Research}, 176(2), 657--690.
  \item Schleyer, M. y Gue, K. (2012). Throughput time distribution analysis for a one-block
        warehouse. \emph{Transportation Research Part E}, 48, 652--666.
  \item Valle, C. A., Beasley, J. E. y da Cunha, A. S. (2017). Optimally solving the joint
        order batching and picker routing problem. \emph{European Journal of Operational
        Research}.
  \item Yang, P., Zhao, Z. y Guo, H. (2020). Order batch picking optimization under different
        storage scenarios for e-commerce warehouses. \emph{Transportation Research Part E},
        136, 101897.
\end{itemize}

Las notas de lectura de estos trabajos están en \code{docs/lecturas-2026-10.md}.

\end{document}
